\documentclass[11pt]{article}

    \usepackage[breakable]{tcolorbox}
    \usepackage{parskip} % Stop auto-indenting (to mimic markdown behaviour)
    

    % Basic figure setup, for now with no caption control since it's done
    % automatically by Pandoc (which extracts ![](path) syntax from Markdown).
    \usepackage{graphicx}
    % Keep aspect ratio if custom image width or height is specified
    \setkeys{Gin}{keepaspectratio}
    % Maintain compatibility with old templates. Remove in nbconvert 6.0
    \let\Oldincludegraphics\includegraphics
    % Ensure that by default, figures have no caption (until we provide a
    % proper Figure object with a Caption API and a way to capture that
    % in the conversion process - todo).
    \usepackage{caption}
    \DeclareCaptionFormat{nocaption}{}
    \captionsetup{format=nocaption,aboveskip=0pt,belowskip=0pt}

    \usepackage{float}
    \floatplacement{figure}{H} % forces figures to be placed at the correct location
    \usepackage{xcolor} % Allow colors to be defined
    \usepackage{enumerate} % Needed for markdown enumerations to work
    \usepackage{geometry} % Used to adjust the document margins
    \usepackage{amsmath} % Equations
    \usepackage{amssymb} % Equations
    \usepackage{textcomp} % defines textquotesingle
    % Hack from http://tex.stackexchange.com/a/47451/13684:
    \AtBeginDocument{%
        \def\PYZsq{\textquotesingle}% Upright quotes in Pygmentized code
    }
    \usepackage{upquote} % Upright quotes for verbatim code
    \usepackage{eurosym} % defines \euro

    \usepackage{iftex}
    \ifPDFTeX
        \usepackage[T1]{fontenc}
        \IfFileExists{alphabeta.sty}{
              \usepackage{alphabeta}
          }{
              \usepackage[mathletters]{ucs}
              \usepackage[utf8x]{inputenc}
          }
    \else
        \usepackage{fontspec}
        \usepackage{unicode-math}
    \fi

    \usepackage{fancyvrb} % verbatim replacement that allows latex
    \usepackage{grffile} % extends the file name processing of package graphics
                         % to support a larger range
    \makeatletter % fix for old versions of grffile with XeLaTeX
    \@ifpackagelater{grffile}{2019/11/01}
    {
      % Do nothing on new versions
    }
    {
      \def\Gread@@xetex#1{%
        \IfFileExists{"\Gin@base".bb}%
        {\Gread@eps{\Gin@base.bb}}%
        {\Gread@@xetex@aux#1}%
      }
    }
    \makeatother
    \usepackage[Export]{adjustbox} % Used to constrain images to a maximum size
    \adjustboxset{max size={0.9\linewidth}{0.9\paperheight}}

    % The hyperref package gives us a pdf with properly built
    % internal navigation ('pdf bookmarks' for the table of contents,
    % internal cross-reference links, web links for URLs, etc.)
    \usepackage{hyperref}
    % The default LaTeX title has an obnoxious amount of whitespace. By default,
    % titling removes some of it. It also provides customization options.
    \usepackage{titling}
    \usepackage{longtable} % longtable support required by pandoc >1.10
    \usepackage{booktabs}  % table support for pandoc > 1.12.2
    \usepackage{array}     % table support for pandoc >= 2.11.3
    \usepackage{calc}      % table minipage width calculation for pandoc >= 2.11.1
    \usepackage[inline]{enumitem} % IRkernel/repr support (it uses the enumerate* environment)
    \usepackage[normalem]{ulem} % ulem is needed to support strikethroughs (\sout)
                                % normalem makes italics be italics, not underlines
    \usepackage{soul}      % strikethrough (\st) support for pandoc >= 3.0.0
    \usepackage{mathrsfs}
    

    
    % Colors for the hyperref package
    \definecolor{urlcolor}{rgb}{0,.145,.698}
    \definecolor{linkcolor}{rgb}{.71,0.21,0.01}
    \definecolor{citecolor}{rgb}{.12,.54,.11}

    % ANSI colors
    \definecolor{ansi-black}{HTML}{3E424D}
    \definecolor{ansi-black-intense}{HTML}{282C36}
    \definecolor{ansi-red}{HTML}{E75C58}
    \definecolor{ansi-red-intense}{HTML}{B22B31}
    \definecolor{ansi-green}{HTML}{00A250}
    \definecolor{ansi-green-intense}{HTML}{007427}
    \definecolor{ansi-yellow}{HTML}{DDB62B}
    \definecolor{ansi-yellow-intense}{HTML}{B27D12}
    \definecolor{ansi-blue}{HTML}{208FFB}
    \definecolor{ansi-blue-intense}{HTML}{0065CA}
    \definecolor{ansi-magenta}{HTML}{D160C4}
    \definecolor{ansi-magenta-intense}{HTML}{A03196}
    \definecolor{ansi-cyan}{HTML}{60C6C8}
    \definecolor{ansi-cyan-intense}{HTML}{258F8F}
    \definecolor{ansi-white}{HTML}{C5C1B4}
    \definecolor{ansi-white-intense}{HTML}{A1A6B2}
    \definecolor{ansi-default-inverse-fg}{HTML}{FFFFFF}
    \definecolor{ansi-default-inverse-bg}{HTML}{000000}

    % common color for the border for error outputs.
    \definecolor{outerrorbackground}{HTML}{FFDFDF}

    % commands and environments needed by pandoc snippets
    % extracted from the output of `pandoc -s`
    \providecommand{\tightlist}{%
      \setlength{\itemsep}{0pt}\setlength{\parskip}{0pt}}
    \DefineVerbatimEnvironment{Highlighting}{Verbatim}{commandchars=\\\{\}}
    % Add ',fontsize=\small' for more characters per line
    \newenvironment{Shaded}{}{}
    \newcommand{\KeywordTok}[1]{\textcolor[rgb]{0.00,0.44,0.13}{\textbf{{#1}}}}
    \newcommand{\DataTypeTok}[1]{\textcolor[rgb]{0.56,0.13,0.00}{{#1}}}
    \newcommand{\DecValTok}[1]{\textcolor[rgb]{0.25,0.63,0.44}{{#1}}}
    \newcommand{\BaseNTok}[1]{\textcolor[rgb]{0.25,0.63,0.44}{{#1}}}
    \newcommand{\FloatTok}[1]{\textcolor[rgb]{0.25,0.63,0.44}{{#1}}}
    \newcommand{\CharTok}[1]{\textcolor[rgb]{0.25,0.44,0.63}{{#1}}}
    \newcommand{\StringTok}[1]{\textcolor[rgb]{0.25,0.44,0.63}{{#1}}}
    \newcommand{\CommentTok}[1]{\textcolor[rgb]{0.38,0.63,0.69}{\textit{{#1}}}}
    \newcommand{\OtherTok}[1]{\textcolor[rgb]{0.00,0.44,0.13}{{#1}}}
    \newcommand{\AlertTok}[1]{\textcolor[rgb]{1.00,0.00,0.00}{\textbf{{#1}}}}
    \newcommand{\FunctionTok}[1]{\textcolor[rgb]{0.02,0.16,0.49}{{#1}}}
    \newcommand{\RegionMarkerTok}[1]{{#1}}
    \newcommand{\ErrorTok}[1]{\textcolor[rgb]{1.00,0.00,0.00}{\textbf{{#1}}}}
    \newcommand{\NormalTok}[1]{{#1}}

    % Additional commands for more recent versions of Pandoc
    \newcommand{\ConstantTok}[1]{\textcolor[rgb]{0.53,0.00,0.00}{{#1}}}
    \newcommand{\SpecialCharTok}[1]{\textcolor[rgb]{0.25,0.44,0.63}{{#1}}}
    \newcommand{\VerbatimStringTok}[1]{\textcolor[rgb]{0.25,0.44,0.63}{{#1}}}
    \newcommand{\SpecialStringTok}[1]{\textcolor[rgb]{0.73,0.40,0.53}{{#1}}}
    \newcommand{\ImportTok}[1]{{#1}}
    \newcommand{\DocumentationTok}[1]{\textcolor[rgb]{0.73,0.13,0.13}{\textit{{#1}}}}
    \newcommand{\AnnotationTok}[1]{\textcolor[rgb]{0.38,0.63,0.69}{\textbf{\textit{{#1}}}}}
    \newcommand{\CommentVarTok}[1]{\textcolor[rgb]{0.38,0.63,0.69}{\textbf{\textit{{#1}}}}}
    \newcommand{\VariableTok}[1]{\textcolor[rgb]{0.10,0.09,0.49}{{#1}}}
    \newcommand{\ControlFlowTok}[1]{\textcolor[rgb]{0.00,0.44,0.13}{\textbf{{#1}}}}
    \newcommand{\OperatorTok}[1]{\textcolor[rgb]{0.40,0.40,0.40}{{#1}}}
    \newcommand{\BuiltInTok}[1]{{#1}}
    \newcommand{\ExtensionTok}[1]{{#1}}
    \newcommand{\PreprocessorTok}[1]{\textcolor[rgb]{0.74,0.48,0.00}{{#1}}}
    \newcommand{\AttributeTok}[1]{\textcolor[rgb]{0.49,0.56,0.16}{{#1}}}
    \newcommand{\InformationTok}[1]{\textcolor[rgb]{0.38,0.63,0.69}{\textbf{\textit{{#1}}}}}
    \newcommand{\WarningTok}[1]{\textcolor[rgb]{0.38,0.63,0.69}{\textbf{\textit{{#1}}}}}
    \makeatletter
    \newsavebox\pandoc@box
    \newcommand*\pandocbounded[1]{%
      \sbox\pandoc@box{#1}%
      % scaling factors for width and height
      \Gscale@div\@tempa\textheight{\dimexpr\ht\pandoc@box+\dp\pandoc@box\relax}%
      \Gscale@div\@tempb\linewidth{\wd\pandoc@box}%
      % select the smaller of both
      \ifdim\@tempb\p@<\@tempa\p@
        \let\@tempa\@tempb
      \fi
      % scaling accordingly (\@tempa < 1)
      \ifdim\@tempa\p@<\p@
        \scalebox{\@tempa}{\usebox\pandoc@box}%
      % scaling not needed, use as it is
      \else
        \usebox{\pandoc@box}%
      \fi
    }
    \makeatother

    % Define a nice break command that doesn't care if a line doesn't already
    % exist.
    \def\br{\hspace*{\fill} \\* }
    % Math Jax compatibility definitions
    \def\gt{>}
    \def\lt{<}
    \let\Oldtex\TeX
    \let\Oldlatex\LaTeX
    \renewcommand{\TeX}{\textrm{\Oldtex}}
    \renewcommand{\LaTeX}{\textrm{\Oldlatex}}
    % Document parameters
    % Document title
    \title{02-mandat}
    
    
    
    
    
    
    
% Pygments definitions
\makeatletter
\def\PY@reset{\let\PY@it=\relax \let\PY@bf=\relax%
    \let\PY@ul=\relax \let\PY@tc=\relax%
    \let\PY@bc=\relax \let\PY@ff=\relax}
\def\PY@tok#1{\csname PY@tok@#1\endcsname}
\def\PY@toks#1+{\ifx\relax#1\empty\else%
    \PY@tok{#1}\expandafter\PY@toks\fi}
\def\PY@do#1{\PY@bc{\PY@tc{\PY@ul{%
    \PY@it{\PY@bf{\PY@ff{#1}}}}}}}
\def\PY#1#2{\PY@reset\PY@toks#1+\relax+\PY@do{#2}}

\@namedef{PY@tok@w}{\def\PY@tc##1{\textcolor[rgb]{0.73,0.73,0.73}{##1}}}
\@namedef{PY@tok@c}{\let\PY@it=\textit\def\PY@tc##1{\textcolor[rgb]{0.24,0.48,0.48}{##1}}}
\@namedef{PY@tok@cp}{\def\PY@tc##1{\textcolor[rgb]{0.61,0.40,0.00}{##1}}}
\@namedef{PY@tok@k}{\let\PY@bf=\textbf\def\PY@tc##1{\textcolor[rgb]{0.00,0.50,0.00}{##1}}}
\@namedef{PY@tok@kp}{\def\PY@tc##1{\textcolor[rgb]{0.00,0.50,0.00}{##1}}}
\@namedef{PY@tok@kt}{\def\PY@tc##1{\textcolor[rgb]{0.69,0.00,0.25}{##1}}}
\@namedef{PY@tok@o}{\def\PY@tc##1{\textcolor[rgb]{0.40,0.40,0.40}{##1}}}
\@namedef{PY@tok@ow}{\let\PY@bf=\textbf\def\PY@tc##1{\textcolor[rgb]{0.67,0.13,1.00}{##1}}}
\@namedef{PY@tok@nb}{\def\PY@tc##1{\textcolor[rgb]{0.00,0.50,0.00}{##1}}}
\@namedef{PY@tok@nf}{\def\PY@tc##1{\textcolor[rgb]{0.00,0.00,1.00}{##1}}}
\@namedef{PY@tok@nc}{\let\PY@bf=\textbf\def\PY@tc##1{\textcolor[rgb]{0.00,0.00,1.00}{##1}}}
\@namedef{PY@tok@nn}{\let\PY@bf=\textbf\def\PY@tc##1{\textcolor[rgb]{0.00,0.00,1.00}{##1}}}
\@namedef{PY@tok@ne}{\let\PY@bf=\textbf\def\PY@tc##1{\textcolor[rgb]{0.80,0.25,0.22}{##1}}}
\@namedef{PY@tok@nv}{\def\PY@tc##1{\textcolor[rgb]{0.10,0.09,0.49}{##1}}}
\@namedef{PY@tok@no}{\def\PY@tc##1{\textcolor[rgb]{0.53,0.00,0.00}{##1}}}
\@namedef{PY@tok@nl}{\def\PY@tc##1{\textcolor[rgb]{0.46,0.46,0.00}{##1}}}
\@namedef{PY@tok@ni}{\let\PY@bf=\textbf\def\PY@tc##1{\textcolor[rgb]{0.44,0.44,0.44}{##1}}}
\@namedef{PY@tok@na}{\def\PY@tc##1{\textcolor[rgb]{0.41,0.47,0.13}{##1}}}
\@namedef{PY@tok@nt}{\let\PY@bf=\textbf\def\PY@tc##1{\textcolor[rgb]{0.00,0.50,0.00}{##1}}}
\@namedef{PY@tok@nd}{\def\PY@tc##1{\textcolor[rgb]{0.67,0.13,1.00}{##1}}}
\@namedef{PY@tok@s}{\def\PY@tc##1{\textcolor[rgb]{0.73,0.13,0.13}{##1}}}
\@namedef{PY@tok@sd}{\let\PY@it=\textit\def\PY@tc##1{\textcolor[rgb]{0.73,0.13,0.13}{##1}}}
\@namedef{PY@tok@si}{\let\PY@bf=\textbf\def\PY@tc##1{\textcolor[rgb]{0.64,0.35,0.47}{##1}}}
\@namedef{PY@tok@se}{\let\PY@bf=\textbf\def\PY@tc##1{\textcolor[rgb]{0.67,0.36,0.12}{##1}}}
\@namedef{PY@tok@sr}{\def\PY@tc##1{\textcolor[rgb]{0.64,0.35,0.47}{##1}}}
\@namedef{PY@tok@ss}{\def\PY@tc##1{\textcolor[rgb]{0.10,0.09,0.49}{##1}}}
\@namedef{PY@tok@sx}{\def\PY@tc##1{\textcolor[rgb]{0.00,0.50,0.00}{##1}}}
\@namedef{PY@tok@m}{\def\PY@tc##1{\textcolor[rgb]{0.40,0.40,0.40}{##1}}}
\@namedef{PY@tok@gh}{\let\PY@bf=\textbf\def\PY@tc##1{\textcolor[rgb]{0.00,0.00,0.50}{##1}}}
\@namedef{PY@tok@gu}{\let\PY@bf=\textbf\def\PY@tc##1{\textcolor[rgb]{0.50,0.00,0.50}{##1}}}
\@namedef{PY@tok@gd}{\def\PY@tc##1{\textcolor[rgb]{0.63,0.00,0.00}{##1}}}
\@namedef{PY@tok@gi}{\def\PY@tc##1{\textcolor[rgb]{0.00,0.52,0.00}{##1}}}
\@namedef{PY@tok@gr}{\def\PY@tc##1{\textcolor[rgb]{0.89,0.00,0.00}{##1}}}
\@namedef{PY@tok@ge}{\let\PY@it=\textit}
\@namedef{PY@tok@gs}{\let\PY@bf=\textbf}
\@namedef{PY@tok@ges}{\let\PY@bf=\textbf\let\PY@it=\textit}
\@namedef{PY@tok@gp}{\let\PY@bf=\textbf\def\PY@tc##1{\textcolor[rgb]{0.00,0.00,0.50}{##1}}}
\@namedef{PY@tok@go}{\def\PY@tc##1{\textcolor[rgb]{0.44,0.44,0.44}{##1}}}
\@namedef{PY@tok@gt}{\def\PY@tc##1{\textcolor[rgb]{0.00,0.27,0.87}{##1}}}
\@namedef{PY@tok@err}{\def\PY@bc##1{{\setlength{\fboxsep}{\string -\fboxrule}\fcolorbox[rgb]{1.00,0.00,0.00}{1,1,1}{\strut ##1}}}}
\@namedef{PY@tok@kc}{\let\PY@bf=\textbf\def\PY@tc##1{\textcolor[rgb]{0.00,0.50,0.00}{##1}}}
\@namedef{PY@tok@kd}{\let\PY@bf=\textbf\def\PY@tc##1{\textcolor[rgb]{0.00,0.50,0.00}{##1}}}
\@namedef{PY@tok@kn}{\let\PY@bf=\textbf\def\PY@tc##1{\textcolor[rgb]{0.00,0.50,0.00}{##1}}}
\@namedef{PY@tok@kr}{\let\PY@bf=\textbf\def\PY@tc##1{\textcolor[rgb]{0.00,0.50,0.00}{##1}}}
\@namedef{PY@tok@bp}{\def\PY@tc##1{\textcolor[rgb]{0.00,0.50,0.00}{##1}}}
\@namedef{PY@tok@fm}{\def\PY@tc##1{\textcolor[rgb]{0.00,0.00,1.00}{##1}}}
\@namedef{PY@tok@vc}{\def\PY@tc##1{\textcolor[rgb]{0.10,0.09,0.49}{##1}}}
\@namedef{PY@tok@vg}{\def\PY@tc##1{\textcolor[rgb]{0.10,0.09,0.49}{##1}}}
\@namedef{PY@tok@vi}{\def\PY@tc##1{\textcolor[rgb]{0.10,0.09,0.49}{##1}}}
\@namedef{PY@tok@vm}{\def\PY@tc##1{\textcolor[rgb]{0.10,0.09,0.49}{##1}}}
\@namedef{PY@tok@sa}{\def\PY@tc##1{\textcolor[rgb]{0.73,0.13,0.13}{##1}}}
\@namedef{PY@tok@sb}{\def\PY@tc##1{\textcolor[rgb]{0.73,0.13,0.13}{##1}}}
\@namedef{PY@tok@sc}{\def\PY@tc##1{\textcolor[rgb]{0.73,0.13,0.13}{##1}}}
\@namedef{PY@tok@dl}{\def\PY@tc##1{\textcolor[rgb]{0.73,0.13,0.13}{##1}}}
\@namedef{PY@tok@s2}{\def\PY@tc##1{\textcolor[rgb]{0.73,0.13,0.13}{##1}}}
\@namedef{PY@tok@sh}{\def\PY@tc##1{\textcolor[rgb]{0.73,0.13,0.13}{##1}}}
\@namedef{PY@tok@s1}{\def\PY@tc##1{\textcolor[rgb]{0.73,0.13,0.13}{##1}}}
\@namedef{PY@tok@mb}{\def\PY@tc##1{\textcolor[rgb]{0.40,0.40,0.40}{##1}}}
\@namedef{PY@tok@mf}{\def\PY@tc##1{\textcolor[rgb]{0.40,0.40,0.40}{##1}}}
\@namedef{PY@tok@mh}{\def\PY@tc##1{\textcolor[rgb]{0.40,0.40,0.40}{##1}}}
\@namedef{PY@tok@mi}{\def\PY@tc##1{\textcolor[rgb]{0.40,0.40,0.40}{##1}}}
\@namedef{PY@tok@il}{\def\PY@tc##1{\textcolor[rgb]{0.40,0.40,0.40}{##1}}}
\@namedef{PY@tok@mo}{\def\PY@tc##1{\textcolor[rgb]{0.40,0.40,0.40}{##1}}}
\@namedef{PY@tok@ch}{\let\PY@it=\textit\def\PY@tc##1{\textcolor[rgb]{0.24,0.48,0.48}{##1}}}
\@namedef{PY@tok@cm}{\let\PY@it=\textit\def\PY@tc##1{\textcolor[rgb]{0.24,0.48,0.48}{##1}}}
\@namedef{PY@tok@cpf}{\let\PY@it=\textit\def\PY@tc##1{\textcolor[rgb]{0.24,0.48,0.48}{##1}}}
\@namedef{PY@tok@c1}{\let\PY@it=\textit\def\PY@tc##1{\textcolor[rgb]{0.24,0.48,0.48}{##1}}}
\@namedef{PY@tok@cs}{\let\PY@it=\textit\def\PY@tc##1{\textcolor[rgb]{0.24,0.48,0.48}{##1}}}

\def\PYZbs{\char`\\}
\def\PYZus{\char`\_}
\def\PYZob{\char`\{}
\def\PYZcb{\char`\}}
\def\PYZca{\char`\^}
\def\PYZam{\char`\&}
\def\PYZlt{\char`\<}
\def\PYZgt{\char`\>}
\def\PYZsh{\char`\#}
\def\PYZpc{\char`\%}
\def\PYZdl{\char`\$}
\def\PYZhy{\char`\-}
\def\PYZsq{\char`\'}
\def\PYZdq{\char`\"}
\def\PYZti{\char`\~}
% for compatibility with earlier versions
\def\PYZat{@}
\def\PYZlb{[}
\def\PYZrb{]}
\makeatother


    % For linebreaks inside Verbatim environment from package fancyvrb.
    \makeatletter
        \newbox\Wrappedcontinuationbox
        \newbox\Wrappedvisiblespacebox
        \newcommand*\Wrappedvisiblespace {\textcolor{red}{\textvisiblespace}}
        \newcommand*\Wrappedcontinuationsymbol {\textcolor{red}{\llap{\tiny$\m@th\hookrightarrow$}}}
        \newcommand*\Wrappedcontinuationindent {3ex }
        \newcommand*\Wrappedafterbreak {\kern\Wrappedcontinuationindent\copy\Wrappedcontinuationbox}
        % Take advantage of the already applied Pygments mark-up to insert
        % potential linebreaks for TeX processing.
        %        {, <, #, %, $, ' and ": go to next line.
        %        _, }, ^, &, >, - and ~: stay at end of broken line.
        % Use of \textquotesingle for straight quote.
        \newcommand*\Wrappedbreaksatspecials {%
            \def\PYGZus{\discretionary{\char`\_}{\Wrappedafterbreak}{\char`\_}}%
            \def\PYGZob{\discretionary{}{\Wrappedafterbreak\char`\{}{\char`\{}}%
            \def\PYGZcb{\discretionary{\char`\}}{\Wrappedafterbreak}{\char`\}}}%
            \def\PYGZca{\discretionary{\char`\^}{\Wrappedafterbreak}{\char`\^}}%
            \def\PYGZam{\discretionary{\char`\&}{\Wrappedafterbreak}{\char`\&}}%
            \def\PYGZlt{\discretionary{}{\Wrappedafterbreak\char`\<}{\char`\<}}%
            \def\PYGZgt{\discretionary{\char`\>}{\Wrappedafterbreak}{\char`\>}}%
            \def\PYGZsh{\discretionary{}{\Wrappedafterbreak\char`\#}{\char`\#}}%
            \def\PYGZpc{\discretionary{}{\Wrappedafterbreak\char`\%}{\char`\%}}%
            \def\PYGZdl{\discretionary{}{\Wrappedafterbreak\char`\$}{\char`\$}}%
            \def\PYGZhy{\discretionary{\char`\-}{\Wrappedafterbreak}{\char`\-}}%
            \def\PYGZsq{\discretionary{}{\Wrappedafterbreak\textquotesingle}{\textquotesingle}}%
            \def\PYGZdq{\discretionary{}{\Wrappedafterbreak\char`\"}{\char`\"}}%
            \def\PYGZti{\discretionary{\char`\~}{\Wrappedafterbreak}{\char`\~}}%
        }
        % Some characters . , ; ? ! / are not pygmentized.
        % This macro makes them "active" and they will insert potential linebreaks
        \newcommand*\Wrappedbreaksatpunct {%
            \lccode`\~`\.\lowercase{\def~}{\discretionary{\hbox{\char`\.}}{\Wrappedafterbreak}{\hbox{\char`\.}}}%
            \lccode`\~`\,\lowercase{\def~}{\discretionary{\hbox{\char`\,}}{\Wrappedafterbreak}{\hbox{\char`\,}}}%
            \lccode`\~`\;\lowercase{\def~}{\discretionary{\hbox{\char`\;}}{\Wrappedafterbreak}{\hbox{\char`\;}}}%
            \lccode`\~`\:\lowercase{\def~}{\discretionary{\hbox{\char`\:}}{\Wrappedafterbreak}{\hbox{\char`\:}}}%
            \lccode`\~`\?\lowercase{\def~}{\discretionary{\hbox{\char`\?}}{\Wrappedafterbreak}{\hbox{\char`\?}}}%
            \lccode`\~`\!\lowercase{\def~}{\discretionary{\hbox{\char`\!}}{\Wrappedafterbreak}{\hbox{\char`\!}}}%
            \lccode`\~`\/\lowercase{\def~}{\discretionary{\hbox{\char`\/}}{\Wrappedafterbreak}{\hbox{\char`\/}}}%
            \catcode`\.\active
            \catcode`\,\active
            \catcode`\;\active
            \catcode`\:\active
            \catcode`\?\active
            \catcode`\!\active
            \catcode`\/\active
            \lccode`\~`\~
        }
    \makeatother

    \let\OriginalVerbatim=\Verbatim
    \makeatletter
    \renewcommand{\Verbatim}[1][1]{%
        %\parskip\z@skip
        \sbox\Wrappedcontinuationbox {\Wrappedcontinuationsymbol}%
        \sbox\Wrappedvisiblespacebox {\FV@SetupFont\Wrappedvisiblespace}%
        \def\FancyVerbFormatLine ##1{\hsize\linewidth
            \vtop{\raggedright\hyphenpenalty\z@\exhyphenpenalty\z@
                \doublehyphendemerits\z@\finalhyphendemerits\z@
                \strut ##1\strut}%
        }%
        % If the linebreak is at a space, the latter will be displayed as visible
        % space at end of first line, and a continuation symbol starts next line.
        % Stretch/shrink are however usually zero for typewriter font.
        \def\FV@Space {%
            \nobreak\hskip\z@ plus\fontdimen3\font minus\fontdimen4\font
            \discretionary{\copy\Wrappedvisiblespacebox}{\Wrappedafterbreak}
            {\kern\fontdimen2\font}%
        }%

        % Allow breaks at special characters using \PYG... macros.
        \Wrappedbreaksatspecials
        % Breaks at punctuation characters . , ; ? ! and / need catcode=\active
        \OriginalVerbatim[#1,codes*=\Wrappedbreaksatpunct]%
    }
    \makeatother

    % Exact colors from NB
    \definecolor{incolor}{HTML}{303F9F}
    \definecolor{outcolor}{HTML}{D84315}
    \definecolor{cellborder}{HTML}{CFCFCF}
    \definecolor{cellbackground}{HTML}{F7F7F7}

    % prompt
    \makeatletter
    \newcommand{\boxspacing}{\kern\kvtcb@left@rule\kern\kvtcb@boxsep}
    \makeatother
    \newcommand{\prompt}[4]{
        {\ttfamily\llap{{\color{#2}[#3]:\hspace{3pt}#4}}\vspace{-\baselineskip}}
    }
    

    
    % Prevent overflowing lines due to hard-to-break entities
    \sloppy
    % Setup hyperref package
    \hypersetup{
      breaklinks=true,  % so long urls are correctly broken across lines
      colorlinks=true,
      urlcolor=urlcolor,
      linkcolor=linkcolor,
      citecolor=citecolor,
      }
    % Slightly bigger margins than the latex defaults
    
    \geometry{verbose,tmargin=1in,bmargin=1in,lmargin=1in,rmargin=1in}
    
    

\begin{document}
    
    \maketitle
    
    

    
    \begin{tcolorbox}[breakable, size=fbox, boxrule=1pt, pad at break*=1mm,colback=cellbackground, colframe=cellborder]
\prompt{In}{incolor}{20}{\boxspacing}
\begin{Verbatim}[commandchars=\\\{\}]
\PY{k+kn}{import}\PY{+w}{ }\PY{n+nn}{marimo}\PY{+w}{ }\PY{k}{as}\PY{+w}{ }\PY{n+nn}{mo}
\PY{k+kn}{import}\PY{+w}{ }\PY{n+nn}{numpy}\PY{+w}{ }\PY{k}{as}\PY{+w}{ }\PY{n+nn}{np}
\PY{k+kn}{import}\PY{+w}{ }\PY{n+nn}{polars}\PY{+w}{ }\PY{k}{as}\PY{+w}{ }\PY{n+nn}{pl}
\PY{k+kn}{from}\PY{+w}{ }\PY{n+nn}{scipy}\PY{+w}{ }\PY{k+kn}{import} \PY{n}{stats}
\PY{k+kn}{import}\PY{+w}{ }\PY{n+nn}{matplotlib}\PY{n+nn}{.}\PY{n+nn}{pyplot}\PY{+w}{ }\PY{k}{as}\PY{+w}{ }\PY{n+nn}{plt}
\PY{k+kn}{import}\PY{+w}{ }\PY{n+nn}{matplotlib}\PY{+w}{ }\PY{k}{as}\PY{+w}{ }\PY{n+nn}{mpl}
\end{Verbatim}
\end{tcolorbox}

    \section{Mandat - 02: Statistiques descriptives et inférence
statistique}\label{mandat---02-statistiques-descriptives-et-infuxe9rence-statistique}

    \begin{tcolorbox}[breakable, size=fbox, boxrule=1pt, pad at break*=1mm,colback=cellbackground, colframe=cellborder]
\prompt{In}{incolor}{21}{\boxspacing}
\begin{Verbatim}[commandchars=\\\{\}]
\PY{n}{dataset} \PY{o}{=} \PY{n}{np}\PY{o}{.}\PY{n}{loadtxt}\PY{p}{(}\PY{l+s+s2}{\PYZdq{}}\PY{l+s+s2}{../input/TempsDeJeu.txt}\PY{l+s+s2}{\PYZdq{}}\PY{p}{)}
\end{Verbatim}
\end{tcolorbox}

    \begin{tcolorbox}[breakable, size=fbox, boxrule=1pt, pad at break*=1mm,colback=cellbackground, colframe=cellborder]
\prompt{In}{incolor}{22}{\boxspacing}
\begin{Verbatim}[commandchars=\\\{\}]
\PY{k}{def}\PY{+w}{ }\PY{n+nf}{hours}\PY{p}{(}\PY{n}{m}\PY{p}{)}\PY{p}{:}
    \PY{k}{return} \PY{n}{m} \PY{o}{/} \PY{l+m+mi}{60}

\PY{n}{stats\PYZus{}dict} \PY{o}{=} \PY{p}{\PYZob{}}
    \PY{l+s+s2}{\PYZdq{}}\PY{l+s+s2}{Moyenne}\PY{l+s+s2}{\PYZdq{}}\PY{p}{:} \PY{n}{dataset}\PY{o}{.}\PY{n}{mean}\PY{p}{(}\PY{p}{)}\PY{p}{,}
    \PY{l+s+s2}{\PYZdq{}}\PY{l+s+s2}{Mediane}\PY{l+s+s2}{\PYZdq{}}\PY{p}{:} \PY{n}{np}\PY{o}{.}\PY{n}{median}\PY{p}{(}\PY{n}{dataset}\PY{p}{)}\PY{p}{,}
    \PY{l+s+s2}{\PYZdq{}}\PY{l+s+s2}{Mode}\PY{l+s+s2}{\PYZdq{}}\PY{p}{:} \PY{n}{stats}\PY{o}{.}\PY{n}{mode}\PY{p}{(}\PY{n}{dataset}\PY{p}{)}\PY{o}{.}\PY{n}{mode}\PY{o}{.}\PY{n}{item}\PY{p}{(}\PY{p}{)}\PY{p}{,}
    \PY{l+s+s2}{\PYZdq{}}\PY{l+s+s2}{Ecart type}\PY{l+s+s2}{\PYZdq{}}\PY{p}{:} \PY{n}{np}\PY{o}{.}\PY{n}{std}\PY{p}{(}\PY{n}{dataset}\PY{p}{,} \PY{n}{ddof}\PY{o}{=}\PY{l+m+mi}{1}\PY{p}{)}\PY{p}{,}
    \PY{l+s+s2}{\PYZdq{}}\PY{l+s+s2}{Variance}\PY{l+s+s2}{\PYZdq{}}\PY{p}{:} \PY{n}{np}\PY{o}{.}\PY{n}{var}\PY{p}{(}\PY{n}{dataset}\PY{p}{,} \PY{n}{ddof}\PY{o}{=}\PY{l+m+mi}{1}\PY{p}{)}\PY{p}{,}
    \PY{l+s+s2}{\PYZdq{}}\PY{l+s+s2}{Min}\PY{l+s+s2}{\PYZdq{}}\PY{p}{:} \PY{n}{dataset}\PY{o}{.}\PY{n}{min}\PY{p}{(}\PY{p}{)}\PY{p}{,}
    \PY{l+s+s2}{\PYZdq{}}\PY{l+s+s2}{Max}\PY{l+s+s2}{\PYZdq{}}\PY{p}{:} \PY{n}{dataset}\PY{o}{.}\PY{n}{max}\PY{p}{(}\PY{p}{)}\PY{p}{,}
    \PY{l+s+s2}{\PYZdq{}}\PY{l+s+s2}{Etendue}\PY{l+s+s2}{\PYZdq{}}\PY{p}{:} \PY{n}{np}\PY{o}{.}\PY{n}{ptp}\PY{p}{(}\PY{n}{dataset}\PY{p}{)}\PY{p}{,}
\PY{p}{\PYZcb{}}

\PY{n}{df} \PY{o}{=} \PY{n}{pl}\PY{o}{.}\PY{n}{DataFrame}\PY{p}{(}
    \PY{p}{\PYZob{}}
        \PY{l+s+s2}{\PYZdq{}}\PY{l+s+s2}{Metrique}\PY{l+s+s2}{\PYZdq{}}\PY{p}{:} \PY{n+nb}{list}\PY{p}{(}\PY{n}{stats\PYZus{}dict}\PY{o}{.}\PY{n}{keys}\PY{p}{(}\PY{p}{)}\PY{p}{)}\PY{p}{,}
        \PY{l+s+s2}{\PYZdq{}}\PY{l+s+s2}{Valeur}\PY{l+s+s2}{\PYZdq{}}\PY{p}{:} \PY{p}{[}\PY{n+nb}{round}\PY{p}{(}\PY{n}{v}\PY{p}{,} \PY{l+m+mi}{2}\PY{p}{)} \PY{k}{for} \PY{n}{v} \PY{o+ow}{in} \PY{n}{stats\PYZus{}dict}\PY{o}{.}\PY{n}{values}\PY{p}{(}\PY{p}{)}\PY{p}{]}\PY{p}{,}
    \PY{p}{\PYZcb{}}
\PY{p}{)}
\end{Verbatim}
\end{tcolorbox}

    \begin{tcolorbox}[breakable, size=fbox, boxrule=1pt, pad at break*=1mm,colback=cellbackground, colframe=cellborder]
\prompt{In}{incolor}{23}{\boxspacing}
\begin{Verbatim}[commandchars=\\\{\}]
\PY{n}{mo}\PY{o}{.}\PY{n}{md}\PY{p}{(}\PY{l+s+sa}{rf}\PY{l+s+s2}{\PYZdq{}\PYZdq{}\PYZdq{}}
\PY{l+s+s2}{\PYZsh{} I) Statistiques descriptives des temps de jeu}

\PY{l+s+s2}{Cette section calcule les statistiques descriptives des temps de jeu (moyenne, médiane, mode, écart\PYZhy{}type, variance, min, max, étendue) conformément à la première exigence du mandat.}

\PY{l+s+si}{\PYZob{}}\PY{n}{mo}\PY{o}{.}\PY{n}{ui}\PY{o}{.}\PY{n}{table}\PY{p}{(}\PY{n}{df}\PY{p}{,}\PY{+w}{ }\PY{n}{selection}\PY{o}{=}\PY{k+kc}{None}\PY{p}{)}\PY{l+s+si}{\PYZcb{}}
\PY{l+s+s2}{\PYZdq{}\PYZdq{}\PYZdq{}}\PY{p}{)}
\end{Verbatim}
\end{tcolorbox}
 
            
\prompt{Out}{outcolor}{23}{}
    
    \section{I) Statistiques descriptives des temps de
jeu}\label{i-statistiques-descriptives-des-temps-de-jeu}

Cette section calcule les statistiques descriptives des temps de jeu
(moyenne, médiane, mode, écart-type, variance, min, max, étendue)
conformément à la première exigence du mandat.

    

    \section{II) Construction de l'histogramme des
données}\label{ii-construction-de-lhistogramme-des-donnuxe9es}

Cette section construit l'histogramme des temps de jeu en utilisant la
loi de Sturges pour déterminer le nombre de classes, puis calcule les
fréquences, fréquences relatives et fréquences cumulées.

    \begin{tcolorbox}[breakable, size=fbox, boxrule=1pt, pad at break*=1mm,colback=cellbackground, colframe=cellborder]
\prompt{In}{incolor}{24}{\boxspacing}
\begin{Verbatim}[commandchars=\\\{\}]
\PY{n}{k} \PY{o}{=} \PY{n+nb}{int}\PY{p}{(}\PY{n}{np}\PY{o}{.}\PY{n}{ceil}\PY{p}{(}\PY{l+m+mi}{1} \PY{o}{+} \PY{n}{np}\PY{o}{.}\PY{n}{log2}\PY{p}{(}\PY{n+nb}{len}\PY{p}{(}\PY{n}{dataset}\PY{p}{)}\PY{p}{)}\PY{p}{)}\PY{p}{)}
\PY{n}{min\PYZus{}val}\PY{p}{,} \PY{n}{max\PYZus{}val} \PY{o}{=} \PY{n}{dataset}\PY{o}{.}\PY{n}{min}\PY{p}{(}\PY{p}{)}\PY{p}{,} \PY{n}{dataset}\PY{o}{.}\PY{n}{max}\PY{p}{(}\PY{p}{)}
\PY{n}{amplitude} \PY{o}{=} \PY{p}{(}\PY{n}{max\PYZus{}val} \PY{o}{\PYZhy{}} \PY{n}{min\PYZus{}val}\PY{p}{)} \PY{o}{/} \PY{n}{k}

\PY{n}{classes} \PY{o}{=} \PY{p}{[}\PY{p}{]}
\PY{n}{limites} \PY{o}{=} \PY{p}{[}\PY{p}{]}
\PY{n}{centres} \PY{o}{=} \PY{p}{[}\PY{p}{]}
\PY{n}{frequences} \PY{o}{=} \PY{p}{[}\PY{p}{]}

\PY{k}{for} \PY{n}{i} \PY{o+ow}{in} \PY{n+nb}{range}\PY{p}{(}\PY{n}{k}\PY{p}{)}\PY{p}{:}
    \PY{n}{lower} \PY{o}{=} \PY{n}{min\PYZus{}val} \PY{o}{+} \PY{n}{i} \PY{o}{*} \PY{n}{amplitude}
    \PY{n}{upper} \PY{o}{=} \PY{n}{min\PYZus{}val} \PY{o}{+} \PY{p}{(}\PY{n}{i} \PY{o}{+} \PY{l+m+mi}{1}\PY{p}{)} \PY{o}{*} \PY{n}{amplitude}
    \PY{n}{centres}\PY{o}{.}\PY{n}{append}\PY{p}{(}\PY{p}{(}\PY{n}{lower} \PY{o}{+} \PY{n}{upper}\PY{p}{)} \PY{o}{/} \PY{l+m+mi}{2}\PY{p}{)}
    \PY{n}{limites}\PY{o}{.}\PY{n}{append}\PY{p}{(}\PY{l+s+sa}{f}\PY{l+s+s2}{\PYZdq{}}\PY{l+s+s2}{[}\PY{l+s+si}{\PYZob{}}\PY{n}{lower}\PY{l+s+si}{:}\PY{l+s+s2}{.0f}\PY{l+s+si}{\PYZcb{}}\PY{l+s+s2}{, }\PY{l+s+si}{\PYZob{}}\PY{n}{upper}\PY{l+s+si}{:}\PY{l+s+s2}{.0f}\PY{l+s+si}{\PYZcb{}}\PY{l+s+s2}{)}\PY{l+s+s2}{\PYZdq{}}\PY{p}{)}

    \PY{k}{if} \PY{n}{i} \PY{o}{==} \PY{n}{k} \PY{o}{\PYZhy{}} \PY{l+m+mi}{1}\PY{p}{:}
        \PY{n}{limites}\PY{p}{[}\PY{o}{\PYZhy{}}\PY{l+m+mi}{1}\PY{p}{]} \PY{o}{=} \PY{l+s+sa}{f}\PY{l+s+s2}{\PYZdq{}}\PY{l+s+s2}{[}\PY{l+s+si}{\PYZob{}}\PY{n}{lower}\PY{l+s+si}{:}\PY{l+s+s2}{.0f}\PY{l+s+si}{\PYZcb{}}\PY{l+s+s2}{, }\PY{l+s+si}{\PYZob{}}\PY{n}{upper}\PY{l+s+si}{:}\PY{l+s+s2}{.0f}\PY{l+s+si}{\PYZcb{}}\PY{l+s+s2}{]}\PY{l+s+s2}{\PYZdq{}}

    \PY{n}{count} \PY{o}{=} \PY{n}{np}\PY{o}{.}\PY{n}{sum}\PY{p}{(}
        \PY{p}{(}\PY{n}{dataset} \PY{o}{\PYZgt{}}\PY{o}{=} \PY{n}{lower}\PY{p}{)} \PY{o}{\PYZam{}} \PY{p}{(}\PY{n}{dataset} \PY{o}{\PYZlt{}} \PY{n}{upper}\PY{p}{)}
        \PY{k}{if} \PY{n}{i} \PY{o}{\PYZlt{}} \PY{n}{k} \PY{o}{\PYZhy{}} \PY{l+m+mi}{1}
        \PY{k}{else} \PY{p}{(}\PY{n}{dataset} \PY{o}{\PYZgt{}}\PY{o}{=} \PY{n}{lower}\PY{p}{)} \PY{o}{\PYZam{}} \PY{p}{(}\PY{n}{dataset} \PY{o}{\PYZlt{}}\PY{o}{=} \PY{n}{upper}\PY{p}{)}
    \PY{p}{)}
    \PY{n}{frequences}\PY{o}{.}\PY{n}{append}\PY{p}{(}\PY{n}{count}\PY{p}{)}

\PY{n}{freq\PYZus{}relatives} \PY{o}{=} \PY{p}{[}\PY{n}{f} \PY{o}{/} \PY{n+nb}{len}\PY{p}{(}\PY{n}{dataset}\PY{p}{)} \PY{k}{for} \PY{n}{f} \PY{o+ow}{in} \PY{n}{frequences}\PY{p}{]}
\PY{n}{freq\PYZus{}cumulees} \PY{o}{=} \PY{n}{np}\PY{o}{.}\PY{n}{cumsum}\PY{p}{(}\PY{n}{frequences}\PY{p}{)}\PY{o}{.}\PY{n}{tolist}\PY{p}{(}\PY{p}{)}

\PY{n}{classes\PYZus{}col} \PY{o}{=} \PY{p}{[}\PY{l+s+sa}{f}\PY{l+s+s2}{\PYZdq{}}\PY{l+s+s2}{Classe }\PY{l+s+si}{\PYZob{}}\PY{n}{i}\PY{+w}{ }\PY{o}{+}\PY{+w}{ }\PY{l+m+mi}{1}\PY{l+s+si}{\PYZcb{}}\PY{l+s+s2}{\PYZdq{}} \PY{k}{for} \PY{n}{i} \PY{o+ow}{in} \PY{n+nb}{range}\PY{p}{(}\PY{n}{k}\PY{p}{)}\PY{p}{]}

\PY{n}{freq\PYZus{}df} \PY{o}{=} \PY{n}{pl}\PY{o}{.}\PY{n}{DataFrame}\PY{p}{(}
    \PY{p}{\PYZob{}}
        \PY{l+s+s2}{\PYZdq{}}\PY{l+s+s2}{Classes}\PY{l+s+s2}{\PYZdq{}}\PY{p}{:} \PY{n}{classes\PYZus{}col}\PY{p}{,}
        \PY{l+s+s2}{\PYZdq{}}\PY{l+s+s2}{Limites}\PY{l+s+s2}{\PYZdq{}}\PY{p}{:} \PY{n}{limites}\PY{p}{,}
        \PY{l+s+s2}{\PYZdq{}}\PY{l+s+s2}{Centres}\PY{l+s+s2}{\PYZdq{}}\PY{p}{:} \PY{p}{[}\PY{n+nb}{round}\PY{p}{(}\PY{n}{c}\PY{p}{,} \PY{l+m+mi}{1}\PY{p}{)} \PY{k}{for} \PY{n}{c} \PY{o+ow}{in} \PY{n}{centres}\PY{p}{]}\PY{p}{,}
        \PY{l+s+s2}{\PYZdq{}}\PY{l+s+s2}{Freq. relatives}\PY{l+s+s2}{\PYZdq{}}\PY{p}{:} \PY{p}{[}\PY{n+nb}{round}\PY{p}{(}\PY{n}{f}\PY{p}{,} \PY{l+m+mi}{4}\PY{p}{)} \PY{k}{for} \PY{n}{f} \PY{o+ow}{in} \PY{n}{freq\PYZus{}relatives}\PY{p}{]}\PY{p}{,}
        \PY{l+s+s2}{\PYZdq{}}\PY{l+s+s2}{Freq. cumulees}\PY{l+s+s2}{\PYZdq{}}\PY{p}{:} \PY{n}{freq\PYZus{}cumulees}\PY{p}{,}
    \PY{p}{\PYZcb{}}
\PY{p}{)}
\end{Verbatim}
\end{tcolorbox}

    \begin{tcolorbox}[breakable, size=fbox, boxrule=1pt, pad at break*=1mm,colback=cellbackground, colframe=cellborder]
\prompt{In}{incolor}{25}{\boxspacing}
\begin{Verbatim}[commandchars=\\\{\}]
\PY{n}{mo}\PY{o}{.}\PY{n}{md}\PY{p}{(}\PY{l+s+sa}{rf}\PY{l+s+s2}{\PYZdq{}\PYZdq{}\PYZdq{}}
\PY{l+s+s2}{\PYZsh{}\PYZsh{} Nombre de classes}
\PY{l+s+s2}{La loi de Sturges est utilisée pour déterminer le nombre de classes:}

\PY{l+s+s2}{\PYZdl{}\PYZdl{}}
\PY{l+s+s2}{k = }\PY{l+s+s2}{\PYZbs{}}\PY{l+s+s2}{lceil 1 + }\PY{l+s+s2}{\PYZbs{}}\PY{l+s+s2}{log\PYZus{}2(n) }\PY{l+s+s2}{\PYZbs{}}\PY{l+s+s2}{rceil = }\PY{l+s+si}{\PYZob{}}\PY{n}{k}\PY{l+s+si}{\PYZcb{}}
\PY{l+s+s2}{\PYZdl{}\PYZdl{}}

\PY{l+s+s2}{Où:}

\PY{l+s+s2}{| | |}
\PY{l+s+s2}{|\PYZhy{}|\PYZhy{}|}
\PY{l+s+s2}{| k | nombre de classes |}
\PY{l+s+s2}{| n | taille de l}\PY{l+s+s2}{\PYZsq{}}\PY{l+s+s2}{échantillon |}

\PY{l+s+s2}{\PYZsh{}\PYZsh{} Population du tableau de valeurs}

\PY{l+s+si}{\PYZob{}}\PY{n}{mo}\PY{o}{.}\PY{n}{ui}\PY{o}{.}\PY{n}{table}\PY{p}{(}\PY{n}{freq\PYZus{}df}\PY{p}{,}\PY{+w}{ }\PY{n}{selection}\PY{o}{=}\PY{k+kc}{None}\PY{p}{)}\PY{l+s+si}{\PYZcb{}}
\PY{l+s+s2}{\PYZdq{}\PYZdq{}\PYZdq{}}\PY{p}{)}
\end{Verbatim}
\end{tcolorbox}
 
            
\prompt{Out}{outcolor}{25}{}
    
    \subsection{Nombre de classes}\label{nombre-de-classes}

La loi de Sturges est utilisée pour déterminer le nombre de classes:

\[
k = \lceil 1 + \log_2(n) \rceil = 8
\]

Où:

\begin{longtable}[]{@{}ll@{}}
\toprule\noalign{}
\endhead
\bottomrule\noalign{}
\endlastfoot
k & nombre de classes \\
n & taille de l'échantillon \\
\end{longtable}

\subsection{Population du tableau de
valeurs}\label{population-du-tableau-de-valeurs}

    

    \begin{tcolorbox}[breakable, size=fbox, boxrule=1pt, pad at break*=1mm,colback=cellbackground, colframe=cellborder]
\prompt{In}{incolor}{26}{\boxspacing}
\begin{Verbatim}[commandchars=\\\{\}]
\PY{n}{plt}\PY{o}{.}\PY{n}{figure}\PY{p}{(}\PY{n}{figsize}\PY{o}{=}\PY{p}{(}\PY{l+m+mi}{10}\PY{p}{,} \PY{l+m+mi}{6}\PY{p}{)}\PY{p}{)}
\PY{n}{plt}\PY{o}{.}\PY{n}{hist}\PY{p}{(}\PY{n}{dataset}\PY{p}{,} \PY{n}{bins}\PY{o}{=}\PY{n}{k}\PY{p}{,} \PY{n}{edgecolor}\PY{o}{=}\PY{l+s+s2}{\PYZdq{}}\PY{l+s+s2}{black}\PY{l+s+s2}{\PYZdq{}}\PY{p}{,} \PY{n}{alpha}\PY{o}{=}\PY{l+m+mf}{0.7}\PY{p}{)}
\PY{n}{plt}\PY{o}{.}\PY{n}{xlabel}\PY{p}{(}\PY{l+s+s2}{\PYZdq{}}\PY{l+s+s2}{Temps de jeu (minutes)}\PY{l+s+s2}{\PYZdq{}}\PY{p}{)}
\PY{n}{plt}\PY{o}{.}\PY{n}{ylabel}\PY{p}{(}\PY{l+s+s2}{\PYZdq{}}\PY{l+s+s2}{Frequence}\PY{l+s+s2}{\PYZdq{}}\PY{p}{)}
\PY{n}{plt}\PY{o}{.}\PY{n}{title}\PY{p}{(}\PY{l+s+sa}{f}\PY{l+s+s2}{\PYZdq{}}\PY{l+s+s2}{Histogramme des temps de jeu (}\PY{l+s+si}{\PYZob{}}\PY{n}{k}\PY{l+s+si}{\PYZcb{}}\PY{l+s+s2}{ classes)}\PY{l+s+s2}{\PYZdq{}}\PY{p}{)}
\PY{n}{plt}\PY{o}{.}\PY{n}{axvline}\PY{p}{(}
    \PY{n}{dataset}\PY{o}{.}\PY{n}{mean}\PY{p}{(}\PY{p}{)}\PY{p}{,}
    \PY{n}{color}\PY{o}{=}\PY{l+s+s2}{\PYZdq{}}\PY{l+s+s2}{red}\PY{l+s+s2}{\PYZdq{}}\PY{p}{,}
    \PY{n}{linestyle}\PY{o}{=}\PY{l+s+s2}{\PYZdq{}}\PY{l+s+s2}{\PYZhy{}\PYZhy{}}\PY{l+s+s2}{\PYZdq{}}\PY{p}{,}
    \PY{n}{label}\PY{o}{=}\PY{l+s+sa}{f}\PY{l+s+s2}{\PYZdq{}}\PY{l+s+s2}{Moyenne: }\PY{l+s+si}{\PYZob{}}\PY{n}{dataset}\PY{o}{.}\PY{n}{mean}\PY{p}{(}\PY{p}{)}\PY{l+s+si}{:}\PY{l+s+s2}{.1f}\PY{l+s+si}{\PYZcb{}}\PY{l+s+s2}{\PYZdq{}}\PY{p}{,}
\PY{p}{)}
\PY{n}{plt}\PY{o}{.}\PY{n}{legend}\PY{p}{(}\PY{p}{)}
\PY{n}{plt}\PY{o}{.}\PY{n}{gca}\PY{p}{(}\PY{p}{)}
\end{Verbatim}
\end{tcolorbox}

            \begin{tcolorbox}[breakable, size=fbox, boxrule=.5pt, pad at break*=1mm, opacityfill=0]
\prompt{Out}{outcolor}{26}{\boxspacing}
\begin{Verbatim}[commandchars=\\\{\}]
<Axes: title=\{'center': 'Histogramme des temps de jeu (8 classes)'\},
xlabel='Temps de jeu (minutes)', ylabel='Frequence'>
\end{Verbatim}
\end{tcolorbox}
        
    \begin{center}
    \adjustimage{max size={0.9\linewidth}{0.9\paperheight}}{output_8_1.png}
    \end{center}
    { \hspace*{\fill} \\}
    
    \section{III) Les données suivent-elles une distribution
normale?}\label{iii-les-donnuxe9es-suivent-elles-une-distribution-normale}

Cette section vérifie si les temps de jeu suivent une distribution
normale en utilisant un test de Khi-deux d'ajustement, comme exigé dans
la troisième partie du mandat. En observant l'histogramme ci-dessus, on
peut voir que les données semblent suivre une distribution
approximativement normale (forme de cloche). Pour le confirmer
quantitativement, nous utilisons le test de Khi-deux.

    \begin{tcolorbox}[breakable, size=fbox, boxrule=1pt, pad at break*=1mm,colback=cellbackground, colframe=cellborder]
\prompt{In}{incolor}{27}{\boxspacing}
\begin{Verbatim}[commandchars=\\\{\}]
\PY{n}{\PYZus{}mean} \PY{o}{=} \PY{n}{dataset}\PY{o}{.}\PY{n}{mean}\PY{p}{(}\PY{p}{)}
\PY{n}{\PYZus{}std} \PY{o}{=} \PY{n}{dataset}\PY{o}{.}\PY{n}{std}\PY{p}{(}\PY{n}{ddof}\PY{o}{=}\PY{l+m+mi}{1}\PY{p}{)}
\PY{n}{\PYZus{}n} \PY{o}{=} \PY{n+nb}{len}\PY{p}{(}\PY{n}{dataset}\PY{p}{)}

\PY{n}{\PYZus{}k\PYZus{}chi2} \PY{o}{=} \PY{n+nb}{int}\PY{p}{(}\PY{n}{np}\PY{o}{.}\PY{n}{ceil}\PY{p}{(}\PY{l+m+mi}{1} \PY{o}{+} \PY{n}{np}\PY{o}{.}\PY{n}{log2}\PY{p}{(}\PY{n}{\PYZus{}n}\PY{p}{)}\PY{p}{)}\PY{p}{)}
\PY{n}{\PYZus{}min\PYZus{}val}\PY{p}{,} \PY{n}{\PYZus{}max\PYZus{}val} \PY{o}{=} \PY{n}{dataset}\PY{o}{.}\PY{n}{min}\PY{p}{(}\PY{p}{)}\PY{p}{,} \PY{n}{dataset}\PY{o}{.}\PY{n}{max}\PY{p}{(}\PY{p}{)}
\PY{n}{\PYZus{}amplitude} \PY{o}{=} \PY{p}{(}\PY{n}{\PYZus{}max\PYZus{}val} \PY{o}{\PYZhy{}} \PY{n}{\PYZus{}min\PYZus{}val}\PY{p}{)} \PY{o}{/} \PY{n}{\PYZus{}k\PYZus{}chi2}

\PY{c+c1}{\PYZsh{} Calculate initial observed and expected frequencies}
\PY{n}{\PYZus{}observed\PYZus{}freq\PYZus{}raw} \PY{o}{=} \PY{p}{[}\PY{p}{]}
\PY{n}{\PYZus{}expected\PYZus{}freq\PYZus{}raw} \PY{o}{=} \PY{p}{[}\PY{p}{]}
\PY{n}{\PYZus{}bins\PYZus{}edges\PYZus{}raw} \PY{o}{=} \PY{p}{[}\PY{p}{]}

\PY{k}{for} \PY{n}{\PYZus{}i} \PY{o+ow}{in} \PY{n+nb}{range}\PY{p}{(}\PY{n}{\PYZus{}k\PYZus{}chi2}\PY{p}{)}\PY{p}{:}
    \PY{n}{\PYZus{}lower} \PY{o}{=} \PY{n}{\PYZus{}min\PYZus{}val} \PY{o}{+} \PY{n}{\PYZus{}i} \PY{o}{*} \PY{n}{\PYZus{}amplitude}
    \PY{n}{\PYZus{}upper} \PY{o}{=} \PY{n}{\PYZus{}min\PYZus{}val} \PY{o}{+} \PY{p}{(}\PY{n}{\PYZus{}i} \PY{o}{+} \PY{l+m+mi}{1}\PY{p}{)} \PY{o}{*} \PY{n}{\PYZus{}amplitude}
    \PY{n}{\PYZus{}bins\PYZus{}edges\PYZus{}raw}\PY{o}{.}\PY{n}{append}\PY{p}{(}\PY{p}{(}\PY{n}{\PYZus{}lower}\PY{p}{,} \PY{n}{\PYZus{}upper}\PY{p}{)}\PY{p}{)}

    \PY{n}{\PYZus{}count\PYZus{}obs} \PY{o}{=} \PY{n}{np}\PY{o}{.}\PY{n}{sum}\PY{p}{(}
        \PY{p}{(}\PY{n}{dataset} \PY{o}{\PYZgt{}}\PY{o}{=} \PY{n}{\PYZus{}lower}\PY{p}{)} \PY{o}{\PYZam{}} \PY{p}{(}\PY{n}{dataset} \PY{o}{\PYZlt{}} \PY{n}{\PYZus{}upper}\PY{p}{)}
        \PY{k}{if} \PY{n}{\PYZus{}i} \PY{o}{\PYZlt{}} \PY{n}{\PYZus{}k\PYZus{}chi2} \PY{o}{\PYZhy{}} \PY{l+m+mi}{1}
        \PY{k}{else} \PY{p}{(}\PY{n}{dataset} \PY{o}{\PYZgt{}}\PY{o}{=} \PY{n}{\PYZus{}lower}\PY{p}{)} \PY{o}{\PYZam{}} \PY{p}{(}\PY{n}{dataset} \PY{o}{\PYZlt{}}\PY{o}{=} \PY{n}{\PYZus{}upper}\PY{p}{)}
    \PY{p}{)}
    \PY{n}{\PYZus{}observed\PYZus{}freq\PYZus{}raw}\PY{o}{.}\PY{n}{append}\PY{p}{(}\PY{n}{\PYZus{}count\PYZus{}obs}\PY{p}{)}

    \PY{n}{\PYZus{}prob\PYZus{}lower} \PY{o}{=} \PY{n}{stats}\PY{o}{.}\PY{n}{norm}\PY{o}{.}\PY{n}{cdf}\PY{p}{(}\PY{n}{\PYZus{}lower}\PY{p}{,} \PY{n}{loc}\PY{o}{=}\PY{n}{\PYZus{}mean}\PY{p}{,} \PY{n}{scale}\PY{o}{=}\PY{n}{\PYZus{}std}\PY{p}{)}
    \PY{n}{\PYZus{}prob\PYZus{}upper} \PY{o}{=} \PY{n}{stats}\PY{o}{.}\PY{n}{norm}\PY{o}{.}\PY{n}{cdf}\PY{p}{(}\PY{n}{\PYZus{}upper}\PY{p}{,} \PY{n}{loc}\PY{o}{=}\PY{n}{\PYZus{}mean}\PY{p}{,} \PY{n}{scale}\PY{o}{=}\PY{n}{\PYZus{}std}\PY{p}{)}
    \PY{n}{\PYZus{}prob\PYZus{}class} \PY{o}{=} \PY{n}{\PYZus{}prob\PYZus{}upper} \PY{o}{\PYZhy{}} \PY{n}{\PYZus{}prob\PYZus{}lower}
    \PY{n}{\PYZus{}expected\PYZus{}freq\PYZus{}raw}\PY{o}{.}\PY{n}{append}\PY{p}{(}\PY{n}{\PYZus{}n} \PY{o}{*} \PY{n}{\PYZus{}prob\PYZus{}class}\PY{p}{)}

\PY{c+c1}{\PYZsh{} ADJUSTMENT: Merge classes with expected frequency \PYZlt{} 5}
\PY{n}{\PYZus{}observed\PYZus{}freq} \PY{o}{=} \PY{p}{[}\PY{p}{]}
\PY{n}{\PYZus{}expected\PYZus{}freq} \PY{o}{=} \PY{p}{[}\PY{p}{]}
\PY{n}{\PYZus{}bins\PYZus{}edges} \PY{o}{=} \PY{p}{[}\PY{p}{]}

\PY{n}{\PYZus{}i} \PY{o}{=} \PY{l+m+mi}{0}
\PY{k}{while} \PY{n}{\PYZus{}i} \PY{o}{\PYZlt{}} \PY{n+nb}{len}\PY{p}{(}\PY{n}{\PYZus{}expected\PYZus{}freq\PYZus{}raw}\PY{p}{)}\PY{p}{:}
    \PY{c+c1}{\PYZsh{} Start a new class}
    \PY{n}{\PYZus{}obs\PYZus{}merged} \PY{o}{=} \PY{n}{\PYZus{}observed\PYZus{}freq\PYZus{}raw}\PY{p}{[}\PY{n}{\PYZus{}i}\PY{p}{]}
    \PY{n}{\PYZus{}exp\PYZus{}merged} \PY{o}{=} \PY{n}{\PYZus{}expected\PYZus{}freq\PYZus{}raw}\PY{p}{[}\PY{n}{\PYZus{}i}\PY{p}{]}
    \PY{n}{\PYZus{}lower\PYZus{}merged} \PY{o}{=} \PY{n}{\PYZus{}bins\PYZus{}edges\PYZus{}raw}\PY{p}{[}\PY{n}{\PYZus{}i}\PY{p}{]}\PY{p}{[}\PY{l+m+mi}{0}\PY{p}{]}
    \PY{n}{\PYZus{}upper\PYZus{}merged} \PY{o}{=} \PY{n}{\PYZus{}bins\PYZus{}edges\PYZus{}raw}\PY{p}{[}\PY{n}{\PYZus{}i}\PY{p}{]}\PY{p}{[}\PY{l+m+mi}{1}\PY{p}{]}

    \PY{c+c1}{\PYZsh{} Merge with subsequent classes while expected frequency \PYZlt{} 5}
    \PY{k}{while} \PY{n}{\PYZus{}exp\PYZus{}merged} \PY{o}{\PYZlt{}} \PY{l+m+mi}{5} \PY{o+ow}{and} \PY{n}{\PYZus{}i} \PY{o}{+} \PY{l+m+mi}{1} \PY{o}{\PYZlt{}} \PY{n+nb}{len}\PY{p}{(}\PY{n}{\PYZus{}expected\PYZus{}freq\PYZus{}raw}\PY{p}{)}\PY{p}{:}
        \PY{n}{\PYZus{}i} \PY{o}{+}\PY{o}{=} \PY{l+m+mi}{1}
        \PY{n}{\PYZus{}obs\PYZus{}merged} \PY{o}{+}\PY{o}{=} \PY{n}{\PYZus{}observed\PYZus{}freq\PYZus{}raw}\PY{p}{[}\PY{n}{\PYZus{}i}\PY{p}{]}
        \PY{n}{\PYZus{}exp\PYZus{}merged} \PY{o}{+}\PY{o}{=} \PY{n}{\PYZus{}expected\PYZus{}freq\PYZus{}raw}\PY{p}{[}\PY{n}{\PYZus{}i}\PY{p}{]}
        \PY{n}{\PYZus{}upper\PYZus{}merged} \PY{o}{=} \PY{n}{\PYZus{}bins\PYZus{}edges\PYZus{}raw}\PY{p}{[}\PY{n}{\PYZus{}i}\PY{p}{]}\PY{p}{[}\PY{l+m+mi}{1}\PY{p}{]}

    \PY{n}{\PYZus{}observed\PYZus{}freq}\PY{o}{.}\PY{n}{append}\PY{p}{(}\PY{n}{\PYZus{}obs\PYZus{}merged}\PY{p}{)}
    \PY{n}{\PYZus{}expected\PYZus{}freq}\PY{o}{.}\PY{n}{append}\PY{p}{(}\PY{n}{\PYZus{}exp\PYZus{}merged}\PY{p}{)}
    \PY{n}{\PYZus{}bins\PYZus{}edges}\PY{o}{.}\PY{n}{append}\PY{p}{(}\PY{p}{(}\PY{n}{\PYZus{}lower\PYZus{}merged}\PY{p}{,} \PY{n}{\PYZus{}upper\PYZus{}merged}\PY{p}{)}\PY{p}{)}
    \PY{n}{\PYZus{}i} \PY{o}{+}\PY{o}{=} \PY{l+m+mi}{1}

\PY{c+c1}{\PYZsh{} Convert to arrays}
\PY{n}{\PYZus{}observed\PYZus{}freq} \PY{o}{=} \PY{n}{np}\PY{o}{.}\PY{n}{array}\PY{p}{(}\PY{n}{\PYZus{}observed\PYZus{}freq}\PY{p}{)}
\PY{n}{\PYZus{}expected\PYZus{}freq} \PY{o}{=} \PY{n}{np}\PY{o}{.}\PY{n}{array}\PY{p}{(}\PY{n}{\PYZus{}expected\PYZus{}freq}\PY{p}{)}

\PY{n}{chi2\PYZus{}stat} \PY{o}{=} \PY{n}{np}\PY{o}{.}\PY{n}{sum}\PY{p}{(}\PY{p}{(}\PY{n}{\PYZus{}observed\PYZus{}freq} \PY{o}{\PYZhy{}} \PY{n}{\PYZus{}expected\PYZus{}freq}\PY{p}{)} \PY{o}{*}\PY{o}{*} \PY{l+m+mi}{2} \PY{o}{/} \PY{n}{\PYZus{}expected\PYZus{}freq}\PY{p}{)}

\PY{c+c1}{\PYZsh{} Degrees of freedom: (number of classes after merging) \PYZhy{} 1 \PYZhy{} 2}
\PY{n}{ddof\PYZus{}chi2} \PY{o}{=} \PY{n+nb}{len}\PY{p}{(}\PY{n}{\PYZus{}observed\PYZus{}freq}\PY{p}{)} \PY{o}{\PYZhy{}} \PY{l+m+mi}{1} \PY{o}{\PYZhy{}} \PY{l+m+mi}{2}
\PY{n}{p\PYZus{}value} \PY{o}{=} \PY{l+m+mi}{1} \PY{o}{\PYZhy{}} \PY{n}{stats}\PY{o}{.}\PY{n}{chi2}\PY{o}{.}\PY{n}{cdf}\PY{p}{(}\PY{n}{chi2\PYZus{}stat}\PY{p}{,} \PY{n}{df}\PY{o}{=}\PY{n}{ddof\PYZus{}chi2}\PY{p}{)}

\PY{n}{alpha} \PY{o}{=} \PY{l+m+mf}{0.05}
\PY{n}{chi2\PYZus{}critical} \PY{o}{=} \PY{n}{stats}\PY{o}{.}\PY{n}{chi2}\PY{o}{.}\PY{n}{ppf}\PY{p}{(}\PY{l+m+mi}{1} \PY{o}{\PYZhy{}} \PY{n}{alpha}\PY{p}{,} \PY{n}{df}\PY{o}{=}\PY{n}{ddof\PYZus{}chi2}\PY{p}{)}

\PY{n}{chi2\PYZus{}result\PYZus{}df} \PY{o}{=} \PY{n}{pl}\PY{o}{.}\PY{n}{DataFrame}\PY{p}{(}
    \PY{p}{\PYZob{}}
        \PY{l+s+s2}{\PYZdq{}}\PY{l+s+s2}{Classe}\PY{l+s+s2}{\PYZdq{}}\PY{p}{:} \PY{p}{[}
            \PY{l+s+sa}{f}\PY{l+s+s2}{\PYZdq{}}\PY{l+s+s2}{[}\PY{l+s+si}{\PYZob{}}\PY{n}{\PYZus{}lower}\PY{l+s+si}{:}\PY{l+s+s2}{.0f}\PY{l+s+si}{\PYZcb{}}\PY{l+s+s2}{, }\PY{l+s+si}{\PYZob{}}\PY{n}{\PYZus{}upper}\PY{l+s+si}{:}\PY{l+s+s2}{.0f}\PY{l+s+si}{\PYZcb{}}\PY{l+s+s2}{]}\PY{l+s+s2}{\PYZdq{}}
            \PY{k}{for} \PY{n}{\PYZus{}lower}\PY{p}{,} \PY{n}{\PYZus{}upper} \PY{o+ow}{in} \PY{n}{\PYZus{}bins\PYZus{}edges}
        \PY{p}{]}\PY{p}{,}
        \PY{l+s+s2}{\PYZdq{}}\PY{l+s+s2}{Freq. observee}\PY{l+s+s2}{\PYZdq{}}\PY{p}{:} \PY{n}{\PYZus{}observed\PYZus{}freq}\PY{o}{.}\PY{n}{tolist}\PY{p}{(}\PY{p}{)}\PY{p}{,}
        \PY{l+s+s2}{\PYZdq{}}\PY{l+s+s2}{Freq. attendue}\PY{l+s+s2}{\PYZdq{}}\PY{p}{:} \PY{p}{[}\PY{n+nb}{round}\PY{p}{(}\PY{n}{\PYZus{}f}\PY{p}{,} \PY{l+m+mi}{2}\PY{p}{)} \PY{k}{for} \PY{n}{\PYZus{}f} \PY{o+ow}{in} \PY{n}{\PYZus{}expected\PYZus{}freq}\PY{p}{]}\PY{p}{,}
    \PY{p}{\PYZcb{}}
\PY{p}{)}
\end{Verbatim}
\end{tcolorbox}

    \begin{tcolorbox}[breakable, size=fbox, boxrule=1pt, pad at break*=1mm,colback=cellbackground, colframe=cellborder]
\prompt{In}{incolor}{28}{\boxspacing}
\begin{Verbatim}[commandchars=\\\{\}]
\PY{n}{reject\PYZus{}h0} \PY{o}{=} \PY{n}{p\PYZus{}value} \PY{o}{\PYZlt{}} \PY{n}{alpha}
\PY{n}{conclusion} \PY{o}{=} \PY{l+s+s2}{\PYZdq{}}\PY{l+s+s2}{rejeter}\PY{l+s+s2}{\PYZdq{}} \PY{k}{if} \PY{n}{reject\PYZus{}h0} \PY{k}{else} \PY{l+s+s2}{\PYZdq{}}\PY{l+s+s2}{ne pas rejeter}\PY{l+s+s2}{\PYZdq{}}

\PY{n}{mo}\PY{o}{.}\PY{n}{md}\PY{p}{(}\PY{l+s+sa}{rf}\PY{l+s+s2}{\PYZdq{}\PYZdq{}\PYZdq{}}
\PY{l+s+s2}{\PYZsh{}\PYZsh{} Test de Khi\PYZhy{}deux (test de Pearson)}

\PY{l+s+s2}{**Hypothèses:**}
\PY{l+s+s2}{\PYZhy{} H₀: Les données suivent une distribution normale}
\PY{l+s+s2}{\PYZhy{} H₁: Les données ne suivent pas une distribution normale}

\PY{l+s+s2}{**Résultats:**}

\PY{l+s+si}{\PYZob{}}\PY{n}{mo}\PY{o}{.}\PY{n}{ui}\PY{o}{.}\PY{n}{table}\PY{p}{(}\PY{n}{chi2\PYZus{}result\PYZus{}df}\PY{p}{,}\PY{+w}{ }\PY{n}{selection}\PY{o}{=}\PY{k+kc}{None}\PY{p}{)}\PY{l+s+si}{\PYZcb{}}

\PY{l+s+s2}{**Statistiques du test:**}

\PY{l+s+s2}{| Métrique | Valeur |}
\PY{l+s+s2}{|\PYZhy{}\PYZhy{}\PYZhy{}\PYZhy{}\PYZhy{}\PYZhy{}\PYZhy{}\PYZhy{}\PYZhy{}\PYZhy{}|\PYZhy{}\PYZhy{}\PYZhy{}\PYZhy{}\PYZhy{}\PYZhy{}\PYZhy{}\PYZhy{}|}
\PY{l+s+s2}{| Statistique χ² | }\PY{l+s+si}{\PYZob{}}\PY{n}{chi2\PYZus{}stat}\PY{l+s+si}{:}\PY{l+s+s2}{.4f}\PY{l+s+si}{\PYZcb{}}\PY{l+s+s2}{ |}
\PY{l+s+s2}{| Degrés de liberté | }\PY{l+s+si}{\PYZob{}}\PY{n}{ddof\PYZus{}chi2}\PY{l+s+si}{\PYZcb{}}\PY{l+s+s2}{ |}
\PY{l+s+s2}{| Valeur p | }\PY{l+s+si}{\PYZob{}}\PY{n}{p\PYZus{}value}\PY{l+s+si}{:}\PY{l+s+s2}{.4f}\PY{l+s+si}{\PYZcb{}}\PY{l+s+s2}{ |}
\PY{l+s+s2}{| Valeur critique (α=}\PY{l+s+si}{\PYZob{}}\PY{n}{alpha}\PY{l+s+si}{\PYZcb{}}\PY{l+s+s2}{) | }\PY{l+s+si}{\PYZob{}}\PY{n}{chi2\PYZus{}critical}\PY{l+s+si}{:}\PY{l+s+s2}{.4f}\PY{l+s+si}{\PYZcb{}}\PY{l+s+s2}{ |}

\PY{l+s+s2}{**Conclusion:**}

\PY{l+s+s2}{Au seuil de signification α = }\PY{l+s+si}{\PYZob{}}\PY{n}{alpha}\PY{l+s+si}{\PYZcb{}}\PY{l+s+s2}{, on }\PY{l+s+si}{\PYZob{}}\PY{l+s+s2}{\PYZdq{}}\PY{l+s+s2}{**rejette**}\PY{l+s+s2}{\PYZdq{}}\PY{+w}{ }\PY{k}{if}\PY{+w}{ }\PY{n}{reject\PYZus{}h0}\PY{+w}{ }\PY{k}{else}\PY{+w}{ }\PY{l+s+s2}{\PYZdq{}}\PY{l+s+s2}{**ne rejette pas**}\PY{l+s+s2}{\PYZdq{}}\PY{l+s+si}{\PYZcb{}}\PY{l+s+s2}{ l}\PY{l+s+s2}{\PYZsq{}}\PY{l+s+s2}{hypothèse nulle.}

\PY{l+s+si}{\PYZob{}}\PY{l+s+s2}{\PYZdq{}}\PY{l+s+s2}{Les données **ne suivent pas** une distribution normale.}\PY{l+s+s2}{\PYZdq{}}\PY{+w}{ }\PY{k}{if}\PY{+w}{ }\PY{n}{reject\PYZus{}h0}\PY{+w}{ }\PY{k}{else}\PY{+w}{ }\PY{l+s+s2}{\PYZdq{}}\PY{l+s+s2}{Les données **suivent** une distribution normale.}\PY{l+s+s2}{\PYZdq{}}\PY{l+s+si}{\PYZcb{}}
\PY{l+s+s2}{\PYZdq{}\PYZdq{}\PYZdq{}}\PY{p}{)}
\end{Verbatim}
\end{tcolorbox}
 
            
\prompt{Out}{outcolor}{28}{}
    
    \subsection{Test de Khi-deux (test de
Pearson)}\label{test-de-khi-deux-test-de-pearson}

\textbf{Hypothèses:} - H₀: Les données suivent une distribution normale
- H₁: Les données ne suivent pas une distribution normale

\textbf{Résultats:}

\textbf{Statistiques du test:}

\begin{longtable}[]{@{}ll@{}}
\toprule\noalign{}
Métrique & Valeur \\
\midrule\noalign{}
\endhead
\bottomrule\noalign{}
\endlastfoot
Statistique χ² & 3.3721 \\
Degrés de liberté & 4 \\
Valeur p & 0.4976 \\
Valeur critique (α=0.05) & 9.4877 \\
\end{longtable}

\textbf{Conclusion:}

Au seuil de signification α = 0.05, on \textbf{ne rejette pas}
l'hypothèse nulle.

Les données \textbf{suivent} une distribution normale.

    

    \section{IV) Intervalle de confiance pour la
moyenne}\label{iv-intervalle-de-confiance-pour-la-moyenne}

Cette section calcule l'intervalle de confiance à 95\% pour la moyenne
des temps de jeu en utilisant la distribution normale centrée réduite,
comme requis dans la quatrième partie du mandat.

Calculez l'intervalle de confiance pour la moyenne des temps de jeu avec
un niveau de confiance de 95\%. Vous ferez alors usage des tables de la
distribution normale centrée réduite.

    \begin{tcolorbox}[breakable, size=fbox, boxrule=1pt, pad at break*=1mm,colback=cellbackground, colframe=cellborder]
\prompt{In}{incolor}{29}{\boxspacing}
\begin{Verbatim}[commandchars=\\\{\}]
\PY{n}{confidence\PYZus{}level} \PY{o}{=} \PY{l+m+mf}{0.95}
\PY{n}{alpha\PYZus{}ci} \PY{o}{=} \PY{l+m+mi}{1} \PY{o}{\PYZhy{}} \PY{n}{confidence\PYZus{}level}

\PY{n}{sample\PYZus{}mean} \PY{o}{=} \PY{n}{dataset}\PY{o}{.}\PY{n}{mean}\PY{p}{(}\PY{p}{)}
\PY{n}{sample\PYZus{}std} \PY{o}{=} \PY{n}{dataset}\PY{o}{.}\PY{n}{std}\PY{p}{(}\PY{n}{ddof}\PY{o}{=}\PY{l+m+mi}{1}\PY{p}{)}
\PY{n}{n\PYZus{}samples} \PY{o}{=} \PY{n+nb}{len}\PY{p}{(}\PY{n}{dataset}\PY{p}{)}

\PY{n}{z\PYZus{}critical} \PY{o}{=} \PY{n}{stats}\PY{o}{.}\PY{n}{norm}\PY{o}{.}\PY{n}{ppf}\PY{p}{(}\PY{l+m+mi}{1} \PY{o}{\PYZhy{}} \PY{n}{alpha\PYZus{}ci} \PY{o}{/} \PY{l+m+mi}{2}\PY{p}{)}

\PY{n}{standard\PYZus{}error} \PY{o}{=} \PY{n}{sample\PYZus{}std} \PY{o}{/} \PY{n}{np}\PY{o}{.}\PY{n}{sqrt}\PY{p}{(}\PY{n}{n\PYZus{}samples}\PY{p}{)}
\PY{n}{margin\PYZus{}of\PYZus{}error} \PY{o}{=} \PY{n}{z\PYZus{}critical} \PY{o}{*} \PY{n}{standard\PYZus{}error}

\PY{n}{ci\PYZus{}lower} \PY{o}{=} \PY{n}{sample\PYZus{}mean} \PY{o}{\PYZhy{}} \PY{n}{margin\PYZus{}of\PYZus{}error}
\PY{n}{ci\PYZus{}upper} \PY{o}{=} \PY{n}{sample\PYZus{}mean} \PY{o}{+} \PY{n}{margin\PYZus{}of\PYZus{}error}

\PY{n}{ci\PYZus{}df} \PY{o}{=} \PY{n}{pl}\PY{o}{.}\PY{n}{DataFrame}\PY{p}{(}
    \PY{p}{\PYZob{}}
        \PY{l+s+s2}{\PYZdq{}}\PY{l+s+s2}{Paramètre}\PY{l+s+s2}{\PYZdq{}}\PY{p}{:} \PY{p}{[}
            \PY{l+s+s2}{\PYZdq{}}\PY{l+s+s2}{Moyenne de l}\PY{l+s+s2}{\PYZsq{}}\PY{l+s+s2}{échantillon (min)}\PY{l+s+s2}{\PYZdq{}}\PY{p}{,}
            \PY{l+s+s2}{\PYZdq{}}\PY{l+s+s2}{Moyenne de l}\PY{l+s+s2}{\PYZsq{}}\PY{l+s+s2}{échantillon (h)}\PY{l+s+s2}{\PYZdq{}}\PY{p}{,}
            \PY{l+s+s2}{\PYZdq{}}\PY{l+s+s2}{Écart\PYZhy{}type de l}\PY{l+s+s2}{\PYZsq{}}\PY{l+s+s2}{échantillon (min)}\PY{l+s+s2}{\PYZdq{}}\PY{p}{,}
            \PY{l+s+s2}{\PYZdq{}}\PY{l+s+s2}{Taille de l}\PY{l+s+s2}{\PYZsq{}}\PY{l+s+s2}{échantillon}\PY{l+s+s2}{\PYZdq{}}\PY{p}{,}
            \PY{l+s+s2}{\PYZdq{}}\PY{l+s+s2}{Niveau de confiance}\PY{l+s+s2}{\PYZdq{}}\PY{p}{,}
            \PY{l+s+s2}{\PYZdq{}}\PY{l+s+s2}{Valeur critique z}\PY{l+s+s2}{\PYZdq{}}\PY{p}{,}
            \PY{l+s+s2}{\PYZdq{}}\PY{l+s+s2}{Erreur standard (min)}\PY{l+s+s2}{\PYZdq{}}\PY{p}{,}
            \PY{l+s+s2}{\PYZdq{}}\PY{l+s+s2}{Marge d}\PY{l+s+s2}{\PYZsq{}}\PY{l+s+s2}{erreur (min)}\PY{l+s+s2}{\PYZdq{}}\PY{p}{,}
            \PY{l+s+s2}{\PYZdq{}}\PY{l+s+s2}{Limite inférieure (min)}\PY{l+s+s2}{\PYZdq{}}\PY{p}{,}
            \PY{l+s+s2}{\PYZdq{}}\PY{l+s+s2}{Limite supérieure (min)}\PY{l+s+s2}{\PYZdq{}}\PY{p}{,}
            \PY{l+s+s2}{\PYZdq{}}\PY{l+s+s2}{Limite inférieure (h)}\PY{l+s+s2}{\PYZdq{}}\PY{p}{,}
            \PY{l+s+s2}{\PYZdq{}}\PY{l+s+s2}{Limite supérieure (h)}\PY{l+s+s2}{\PYZdq{}}\PY{p}{,}
        \PY{p}{]}\PY{p}{,}
        \PY{l+s+s2}{\PYZdq{}}\PY{l+s+s2}{Valeur}\PY{l+s+s2}{\PYZdq{}}\PY{p}{:} \PY{p}{[}
            \PY{l+s+sa}{f}\PY{l+s+s2}{\PYZdq{}}\PY{l+s+si}{\PYZob{}}\PY{n}{sample\PYZus{}mean}\PY{l+s+si}{:}\PY{l+s+s2}{.2f}\PY{l+s+si}{\PYZcb{}}\PY{l+s+s2}{\PYZdq{}}\PY{p}{,}
            \PY{l+s+sa}{f}\PY{l+s+s2}{\PYZdq{}}\PY{l+s+si}{\PYZob{}}\PY{n}{sample\PYZus{}mean}\PY{l+s+si}{:}\PY{l+s+s2}{.2f}\PY{l+s+si}{\PYZcb{}}\PY{l+s+s2}{\PYZdq{}}\PY{p}{,}
            \PY{l+s+sa}{f}\PY{l+s+s2}{\PYZdq{}}\PY{l+s+si}{\PYZob{}}\PY{n}{sample\PYZus{}std}\PY{l+s+si}{:}\PY{l+s+s2}{.2f}\PY{l+s+si}{\PYZcb{}}\PY{l+s+s2}{\PYZdq{}}\PY{p}{,}
            \PY{l+s+sa}{f}\PY{l+s+s2}{\PYZdq{}}\PY{l+s+si}{\PYZob{}}\PY{n}{n\PYZus{}samples}\PY{l+s+si}{\PYZcb{}}\PY{l+s+s2}{\PYZdq{}}\PY{p}{,}
            \PY{l+s+sa}{f}\PY{l+s+s2}{\PYZdq{}}\PY{l+s+si}{\PYZob{}}\PY{n}{confidence\PYZus{}level}\PY{+w}{ }\PY{o}{*}\PY{+w}{ }\PY{l+m+mi}{100}\PY{l+s+si}{:}\PY{l+s+s2}{.0f}\PY{l+s+si}{\PYZcb{}}\PY{l+s+s2}{\PYZpc{}}\PY{l+s+s2}{\PYZdq{}}\PY{p}{,}
            \PY{l+s+sa}{f}\PY{l+s+s2}{\PYZdq{}}\PY{l+s+si}{\PYZob{}}\PY{n}{z\PYZus{}critical}\PY{l+s+si}{:}\PY{l+s+s2}{.4f}\PY{l+s+si}{\PYZcb{}}\PY{l+s+s2}{\PYZdq{}}\PY{p}{,}
            \PY{l+s+sa}{f}\PY{l+s+s2}{\PYZdq{}}\PY{l+s+si}{\PYZob{}}\PY{n}{standard\PYZus{}error}\PY{l+s+si}{:}\PY{l+s+s2}{.2f}\PY{l+s+si}{\PYZcb{}}\PY{l+s+s2}{\PYZdq{}}\PY{p}{,}
            \PY{l+s+sa}{f}\PY{l+s+s2}{\PYZdq{}}\PY{l+s+si}{\PYZob{}}\PY{n}{margin\PYZus{}of\PYZus{}error}\PY{l+s+si}{:}\PY{l+s+s2}{.2f}\PY{l+s+si}{\PYZcb{}}\PY{l+s+s2}{\PYZdq{}}\PY{p}{,}
            \PY{l+s+sa}{f}\PY{l+s+s2}{\PYZdq{}}\PY{l+s+si}{\PYZob{}}\PY{n}{ci\PYZus{}lower}\PY{l+s+si}{:}\PY{l+s+s2}{.2f}\PY{l+s+si}{\PYZcb{}}\PY{l+s+s2}{\PYZdq{}}\PY{p}{,}
            \PY{l+s+sa}{f}\PY{l+s+s2}{\PYZdq{}}\PY{l+s+si}{\PYZob{}}\PY{n}{ci\PYZus{}upper}\PY{l+s+si}{:}\PY{l+s+s2}{.2f}\PY{l+s+si}{\PYZcb{}}\PY{l+s+s2}{\PYZdq{}}\PY{p}{,}
            \PY{l+s+sa}{f}\PY{l+s+s2}{\PYZdq{}}\PY{l+s+si}{\PYZob{}}\PY{n}{ci\PYZus{}lower}\PY{l+s+si}{:}\PY{l+s+s2}{.2f}\PY{l+s+si}{\PYZcb{}}\PY{l+s+s2}{\PYZdq{}}\PY{p}{,}
            \PY{l+s+sa}{f}\PY{l+s+s2}{\PYZdq{}}\PY{l+s+si}{\PYZob{}}\PY{n}{ci\PYZus{}upper}\PY{l+s+si}{:}\PY{l+s+s2}{.2f}\PY{l+s+si}{\PYZcb{}}\PY{l+s+s2}{\PYZdq{}}\PY{p}{,}
        \PY{p}{]}\PY{p}{,}
    \PY{p}{\PYZcb{}}
\PY{p}{)}
\end{Verbatim}
\end{tcolorbox}

    \begin{tcolorbox}[breakable, size=fbox, boxrule=1pt, pad at break*=1mm,colback=cellbackground, colframe=cellborder]
\prompt{In}{incolor}{30}{\boxspacing}
\begin{Verbatim}[commandchars=\\\{\}]
\PY{n}{mo}\PY{o}{.}\PY{n}{md}\PY{p}{(}\PY{l+s+sa}{rf}\PY{l+s+s2}{\PYZdq{}\PYZdq{}\PYZdq{}}
\PY{l+s+s2}{\PYZsh{}\PYZsh{} Calcul de l}\PY{l+s+s2}{\PYZsq{}}\PY{l+s+s2}{intervalle de confiance}

\PY{l+s+s2}{Pour un niveau de confiance de }\PY{l+s+si}{\PYZob{}}\PY{n}{confidence\PYZus{}level}\PY{+w}{ }\PY{o}{*}\PY{+w}{ }\PY{l+m+mi}{100}\PY{l+s+si}{:}\PY{l+s+s2}{.0f}\PY{l+s+si}{\PYZcb{}}\PY{l+s+s2}{\PYZpc{}, l}\PY{l+s+s2}{\PYZsq{}}\PY{l+s+s2}{intervalle de confiance pour la moyenne est calculé avec la formule:}

\PY{l+s+s2}{\PYZdl{}\PYZdl{}}
\PY{l+s+s2}{\PYZbs{}}\PY{l+s+s2}{bar}\PY{l+s+se}{\PYZob{}\PYZob{}}\PY{l+s+s2}{x}\PY{l+s+se}{\PYZcb{}\PYZcb{}}\PY{l+s+s2}{ }\PY{l+s+s2}{\PYZbs{}}\PY{l+s+s2}{pm z\PYZus{}}\PY{l+s+se}{\PYZob{}\PYZob{}}\PY{l+s+s2}{\PYZbs{}}\PY{l+s+s2}{alpha/2}\PY{l+s+se}{\PYZcb{}\PYZcb{}}\PY{l+s+s2}{ }\PY{l+s+s2}{\PYZbs{}}\PY{l+s+s2}{times }\PY{l+s+s2}{\PYZbs{}}\PY{l+s+s2}{frac}\PY{l+s+se}{\PYZob{}\PYZob{}}\PY{l+s+s2}{s}\PY{l+s+se}{\PYZcb{}\PYZcb{}}\PY{l+s+se}{\PYZob{}\PYZob{}}\PY{l+s+s2}{\PYZbs{}}\PY{l+s+s2}{sqrt}\PY{l+s+se}{\PYZob{}\PYZob{}}\PY{l+s+s2}{n}\PY{l+s+se}{\PYZcb{}\PYZcb{}}\PY{l+s+se}{\PYZcb{}\PYZcb{}}
\PY{l+s+s2}{\PYZdl{}\PYZdl{}}

\PY{l+s+s2}{Où:}

\PY{l+s+s2}{| Symbole | Description |}
\PY{l+s+s2}{|\PYZhy{}\PYZhy{}\PYZhy{}\PYZhy{}\PYZhy{}\PYZhy{}\PYZhy{}\PYZhy{}\PYZhy{}|\PYZhy{}\PYZhy{}\PYZhy{}\PYZhy{}\PYZhy{}\PYZhy{}\PYZhy{}\PYZhy{}\PYZhy{}\PYZhy{}\PYZhy{}\PYZhy{}\PYZhy{}|}
\PY{l+s+s2}{| \PYZdl{}}\PY{l+s+s2}{\PYZbs{}}\PY{l+s+s2}{bar}\PY{l+s+se}{\PYZob{}\PYZob{}}\PY{l+s+s2}{x}\PY{l+s+se}{\PYZcb{}\PYZcb{}}\PY{l+s+s2}{\PYZdl{} | Moyenne de l}\PY{l+s+s2}{\PYZsq{}}\PY{l+s+s2}{échantillon |}
\PY{l+s+s2}{| \PYZdl{}z\PYZus{}}\PY{l+s+se}{\PYZob{}\PYZob{}}\PY{l+s+s2}{\PYZbs{}}\PY{l+s+s2}{alpha/2}\PY{l+s+se}{\PYZcb{}\PYZcb{}}\PY{l+s+s2}{\PYZdl{} | Valeur critique de la distribution normale (pour α=}\PY{l+s+si}{\PYZob{}}\PY{n}{alpha\PYZus{}ci}\PY{l+s+si}{\PYZcb{}}\PY{l+s+s2}{) |}
\PY{l+s+s2}{| \PYZdl{}s\PYZdl{} | Écart\PYZhy{}type de l}\PY{l+s+s2}{\PYZsq{}}\PY{l+s+s2}{échantillon |}
\PY{l+s+s2}{| \PYZdl{}n\PYZdl{} | Taille de l}\PY{l+s+s2}{\PYZsq{}}\PY{l+s+s2}{échantillon |}

\PY{l+s+s2}{**Résultats:**}

\PY{l+s+si}{\PYZob{}}\PY{n}{mo}\PY{o}{.}\PY{n}{ui}\PY{o}{.}\PY{n}{table}\PY{p}{(}\PY{n}{ci\PYZus{}df}\PY{p}{,}\PY{+w}{ }\PY{n}{selection}\PY{o}{=}\PY{k+kc}{None}\PY{p}{)}\PY{l+s+si}{\PYZcb{}}

\PY{l+s+s2}{**Conclusion:**}

\PY{l+s+s2}{Avec un niveau de confiance de }\PY{l+s+si}{\PYZob{}}\PY{n}{confidence\PYZus{}level}\PY{+w}{ }\PY{o}{*}\PY{+w}{ }\PY{l+m+mi}{100}\PY{l+s+si}{:}\PY{l+s+s2}{.0f}\PY{l+s+si}{\PYZcb{}}\PY{l+s+s2}{\PYZpc{}, l}\PY{l+s+s2}{\PYZsq{}}\PY{l+s+s2}{intervalle de confiance pour la moyenne des temps de jeu est:}

\PY{l+s+s2}{\PYZdl{}\PYZdl{}}
\PY{l+s+s2}{[}\PY{l+s+si}{\PYZob{}}\PY{n}{ci\PYZus{}lower}\PY{l+s+si}{:}\PY{l+s+s2}{.2f}\PY{l+s+si}{\PYZcb{}}\PY{l+s+s2}{, }\PY{l+s+si}{\PYZob{}}\PY{n}{ci\PYZus{}upper}\PY{l+s+si}{:}\PY{l+s+s2}{.2f}\PY{l+s+si}{\PYZcb{}}\PY{l+s+s2}{] }\PY{l+s+s2}{\PYZbs{}}\PY{l+s+s2}{text}\PY{l+s+se}{\PYZob{}\PYZob{}}\PY{l+s+s2}{ minutes}\PY{l+s+se}{\PYZcb{}\PYZcb{}}\PY{l+s+s2}{ = [}\PY{l+s+si}{\PYZob{}}\PY{n}{hours}\PY{p}{(}\PY{n}{ci\PYZus{}lower}\PY{p}{)}\PY{l+s+si}{:}\PY{l+s+s2}{.2f}\PY{l+s+si}{\PYZcb{}}\PY{l+s+s2}{, }\PY{l+s+si}{\PYZob{}}\PY{n}{hours}\PY{p}{(}\PY{n}{ci\PYZus{}upper}\PY{p}{)}\PY{l+s+si}{:}\PY{l+s+s2}{.2f}\PY{l+s+si}{\PYZcb{}}\PY{l+s+s2}{] }\PY{l+s+s2}{\PYZbs{}}\PY{l+s+s2}{text}\PY{l+s+se}{\PYZob{}\PYZob{}}\PY{l+s+s2}{ heures}\PY{l+s+se}{\PYZcb{}\PYZcb{}}
\PY{l+s+s2}{\PYZdl{}\PYZdl{}}

\PY{l+s+s2}{Cela signifie que nous sommes }\PY{l+s+si}{\PYZob{}}\PY{n}{confidence\PYZus{}level}\PY{+w}{ }\PY{o}{*}\PY{+w}{ }\PY{l+m+mi}{100}\PY{l+s+si}{:}\PY{l+s+s2}{.0f}\PY{l+s+si}{\PYZcb{}}\PY{l+s+s2}{\PYZpc{} confiants que la vraie moyenne de la population se situe dans cet intervalle.}
\PY{l+s+s2}{\PYZdq{}\PYZdq{}\PYZdq{}}\PY{p}{)}
\end{Verbatim}
\end{tcolorbox}
 
            
\prompt{Out}{outcolor}{30}{}
    
    \subsection{Calcul de l'intervalle de
confiance}\label{calcul-de-lintervalle-de-confiance}

Pour un niveau de confiance de 95\%, l'intervalle de confiance pour la
moyenne est calculé avec la formule:

\[
\bar{x} \pm z_{\alpha/2} \times \frac{s}{\sqrt{n}}
\]

Où:

\begin{longtable}[]{@{}
  >{\raggedright\arraybackslash}p{(\columnwidth - 2\tabcolsep) * \real{0.4091}}
  >{\raggedright\arraybackslash}p{(\columnwidth - 2\tabcolsep) * \real{0.5909}}@{}}
\toprule\noalign{}
\begin{minipage}[b]{\linewidth}\raggedright
Symbole
\end{minipage} & \begin{minipage}[b]{\linewidth}\raggedright
Description
\end{minipage} \\
\midrule\noalign{}
\endhead
\bottomrule\noalign{}
\endlastfoot
\(\bar{x}\) & Moyenne de l'échantillon \\
\(z_{\alpha/2}\) & Valeur critique de la distribution normale (pour
α=0.050000000000000044) \\
\(s\) & Écart-type de l'échantillon \\
\(n\) & Taille de l'échantillon \\
\end{longtable}

\textbf{Résultats:}

\textbf{Conclusion:}

Avec un niveau de confiance de 95\%, l'intervalle de confiance pour la
moyenne des temps de jeu est:

\[
[270.71, 290.45] \text{ minutes} = [4.51, 4.84] \text{ heures}
\]

Cela signifie que nous sommes 95\% confiants que la vraie moyenne de la
population se situe dans cet intervalle.

    

    \section{V) Test d'hypothèse sur la
moyenne}\label{v-test-dhypothuxe8se-sur-la-moyenne}

Cette section effectue un test d'hypothèse unilatéral à gauche sur la
moyenne pour vérifier si le temps de jeu moyen est significativement
inférieur aux 5 heures (300 minutes) estimées par le patron, comme
requis dans la cinquième partie du mandat.

Effectuez un test d'hypothèse approprié sur la moyenne afin d'évaluer si
les données fournissent suffisamment de preuves pour rejeter ou ne pas
rejeter l'hypothèse de votre patron énoncée précédemment. Considérez un
niveau de confiance de 95\%. Quelle est ici l'erreur de première espèce?

    \begin{tcolorbox}[breakable, size=fbox, boxrule=1pt, pad at break*=1mm,colback=cellbackground, colframe=cellborder]
\prompt{In}{incolor}{31}{\boxspacing}
\begin{Verbatim}[commandchars=\\\{\}]
\PY{c+c1}{\PYZsh{} Hypothesis test on the mean}
\PY{c+c1}{\PYZsh{} H0: μ ≥ 300 (boss claims average playtime is at least 300 minutes)}
\PY{c+c1}{\PYZsh{} H1: μ \PYZlt{} 300 (one\PYZhy{}tailed test, left\PYZhy{}tailed)}

\PY{n}{mu\PYZus{}0} \PY{o}{=} \PY{l+m+mi}{300}  \PY{c+c1}{\PYZsh{} hypothesized population mean (5 hours = 300 minutes)}
\PY{n}{alpha\PYZus{}v} \PY{o}{=} \PY{l+m+mf}{0.05}  \PY{c+c1}{\PYZsh{} significance level (95\PYZpc{} confidence)}

\PY{c+c1}{\PYZsh{} Sample statistics}
\PY{n}{x\PYZus{}bar} \PY{o}{=} \PY{n}{dataset}\PY{o}{.}\PY{n}{mean}\PY{p}{(}\PY{p}{)}
\PY{n}{s} \PY{o}{=} \PY{n}{dataset}\PY{o}{.}\PY{n}{std}\PY{p}{(}\PY{n}{ddof}\PY{o}{=}\PY{l+m+mi}{1}\PY{p}{)}
\PY{n}{n} \PY{o}{=} \PY{n+nb}{len}\PY{p}{(}\PY{n}{dataset}\PY{p}{)}

\PY{c+c1}{\PYZsh{} Test statistic (z\PYZhy{}test since n is large)}
\PY{n}{z\PYZus{}stat} \PY{o}{=} \PY{p}{(}\PY{n}{x\PYZus{}bar} \PY{o}{\PYZhy{}} \PY{n}{mu\PYZus{}0}\PY{p}{)} \PY{o}{/} \PY{p}{(}\PY{n}{s} \PY{o}{/} \PY{n}{np}\PY{o}{.}\PY{n}{sqrt}\PY{p}{(}\PY{n}{n}\PY{p}{)}\PY{p}{)}

\PY{c+c1}{\PYZsh{} Critical value for one\PYZhy{}tailed test (left\PYZhy{}tailed)}
\PY{n}{z\PYZus{}critical\PYZus{}v} \PY{o}{=} \PY{n}{stats}\PY{o}{.}\PY{n}{norm}\PY{o}{.}\PY{n}{ppf}\PY{p}{(}\PY{n}{alpha\PYZus{}v}\PY{p}{)}

\PY{c+c1}{\PYZsh{} P\PYZhy{}value for one\PYZhy{}tailed test}
\PY{n}{p\PYZus{}value\PYZus{}v} \PY{o}{=} \PY{n}{stats}\PY{o}{.}\PY{n}{norm}\PY{o}{.}\PY{n}{cdf}\PY{p}{(}\PY{n}{z\PYZus{}stat}\PY{p}{)}

\PY{c+c1}{\PYZsh{} Decision}
\PY{n}{reject\PYZus{}h0\PYZus{}v} \PY{o}{=} \PY{n}{z\PYZus{}stat} \PY{o}{\PYZlt{}} \PY{n}{z\PYZus{}critical\PYZus{}v}

\PY{c+c1}{\PYZsh{} Type I error (alpha) \PYZhy{} probability of rejecting H0 when H0 is true}
\PY{n}{type\PYZus{}i\PYZus{}error} \PY{o}{=} \PY{n}{alpha\PYZus{}v}

\PY{n}{hyp\PYZus{}test\PYZus{}df} \PY{o}{=} \PY{n}{pl}\PY{o}{.}\PY{n}{DataFrame}\PY{p}{(}
    \PY{p}{\PYZob{}}
        \PY{l+s+s2}{\PYZdq{}}\PY{l+s+s2}{Paramètre}\PY{l+s+s2}{\PYZdq{}}\PY{p}{:} \PY{p}{[}
            \PY{l+s+s2}{\PYZdq{}}\PY{l+s+s2}{Hypothèse nulle H₀}\PY{l+s+s2}{\PYZdq{}}\PY{p}{,}
            \PY{l+s+s2}{\PYZdq{}}\PY{l+s+s2}{Hypothèse alternative H₁}\PY{l+s+s2}{\PYZdq{}}\PY{p}{,}
            \PY{l+s+s2}{\PYZdq{}}\PY{l+s+s2}{Moyenne hypothétisée μ₀ (min)}\PY{l+s+s2}{\PYZdq{}}\PY{p}{,}
            \PY{l+s+s2}{\PYZdq{}}\PY{l+s+s2}{Moyenne hypothétisée μ₀ (h)}\PY{l+s+s2}{\PYZdq{}}\PY{p}{,}
            \PY{l+s+s2}{\PYZdq{}}\PY{l+s+s2}{Moyenne échantillonnale x̄ (min)}\PY{l+s+s2}{\PYZdq{}}\PY{p}{,}
            \PY{l+s+s2}{\PYZdq{}}\PY{l+s+s2}{Moyenne échantillonnale x̄ (h)}\PY{l+s+s2}{\PYZdq{}}\PY{p}{,}
            \PY{l+s+s2}{\PYZdq{}}\PY{l+s+s2}{Écart\PYZhy{}type s (min)}\PY{l+s+s2}{\PYZdq{}}\PY{p}{,}
            \PY{l+s+s2}{\PYZdq{}}\PY{l+s+s2}{Taille de l}\PY{l+s+s2}{\PYZsq{}}\PY{l+s+s2}{échantillon n}\PY{l+s+s2}{\PYZdq{}}\PY{p}{,}
            \PY{l+s+s2}{\PYZdq{}}\PY{l+s+s2}{Statistique de test z}\PY{l+s+s2}{\PYZdq{}}\PY{p}{,}
            \PY{l+s+s2}{\PYZdq{}}\PY{l+s+s2}{Valeur critique zα}\PY{l+s+s2}{\PYZdq{}}\PY{p}{,}
            \PY{l+s+s2}{\PYZdq{}}\PY{l+s+s2}{Valeur p}\PY{l+s+s2}{\PYZdq{}}\PY{p}{,}
            \PY{l+s+s2}{\PYZdq{}}\PY{l+s+s2}{Niveau de signification α}\PY{l+s+s2}{\PYZdq{}}\PY{p}{,}
            \PY{l+s+s2}{\PYZdq{}}\PY{l+s+s2}{Erreur de type I (α)}\PY{l+s+s2}{\PYZdq{}}\PY{p}{,}
        \PY{p}{]}\PY{p}{,}
        \PY{l+s+s2}{\PYZdq{}}\PY{l+s+s2}{Valeur}\PY{l+s+s2}{\PYZdq{}}\PY{p}{:} \PY{p}{[}
            \PY{l+s+s2}{\PYZdq{}}\PY{l+s+s2}{μ ≥ 300}\PY{l+s+s2}{\PYZdq{}}\PY{p}{,}
            \PY{l+s+s2}{\PYZdq{}}\PY{l+s+s2}{μ \PYZlt{} 300}\PY{l+s+s2}{\PYZdq{}}\PY{p}{,}
            \PY{l+s+sa}{f}\PY{l+s+s2}{\PYZdq{}}\PY{l+s+si}{\PYZob{}}\PY{n}{mu\PYZus{}0}\PY{l+s+si}{\PYZcb{}}\PY{l+s+s2}{\PYZdq{}}\PY{p}{,}
            \PY{l+s+sa}{f}\PY{l+s+s2}{\PYZdq{}}\PY{l+s+si}{\PYZob{}}\PY{n}{mu\PYZus{}0}\PY{l+s+si}{:}\PY{l+s+s2}{.2f}\PY{l+s+si}{\PYZcb{}}\PY{l+s+s2}{\PYZdq{}}\PY{p}{,}
            \PY{l+s+sa}{f}\PY{l+s+s2}{\PYZdq{}}\PY{l+s+si}{\PYZob{}}\PY{n}{x\PYZus{}bar}\PY{l+s+si}{:}\PY{l+s+s2}{.2f}\PY{l+s+si}{\PYZcb{}}\PY{l+s+s2}{\PYZdq{}}\PY{p}{,}
            \PY{l+s+sa}{f}\PY{l+s+s2}{\PYZdq{}}\PY{l+s+si}{\PYZob{}}\PY{n}{x\PYZus{}bar}\PY{l+s+si}{:}\PY{l+s+s2}{.2f}\PY{l+s+si}{\PYZcb{}}\PY{l+s+s2}{\PYZdq{}}\PY{p}{,}
            \PY{l+s+sa}{f}\PY{l+s+s2}{\PYZdq{}}\PY{l+s+si}{\PYZob{}}\PY{n}{s}\PY{l+s+si}{:}\PY{l+s+s2}{.2f}\PY{l+s+si}{\PYZcb{}}\PY{l+s+s2}{\PYZdq{}}\PY{p}{,}
            \PY{l+s+sa}{f}\PY{l+s+s2}{\PYZdq{}}\PY{l+s+si}{\PYZob{}}\PY{n}{n}\PY{l+s+si}{\PYZcb{}}\PY{l+s+s2}{\PYZdq{}}\PY{p}{,}
            \PY{l+s+sa}{f}\PY{l+s+s2}{\PYZdq{}}\PY{l+s+si}{\PYZob{}}\PY{n}{z\PYZus{}stat}\PY{l+s+si}{:}\PY{l+s+s2}{.4f}\PY{l+s+si}{\PYZcb{}}\PY{l+s+s2}{\PYZdq{}}\PY{p}{,}
            \PY{l+s+sa}{f}\PY{l+s+s2}{\PYZdq{}}\PY{l+s+si}{\PYZob{}}\PY{n}{z\PYZus{}critical\PYZus{}v}\PY{l+s+si}{:}\PY{l+s+s2}{.4f}\PY{l+s+si}{\PYZcb{}}\PY{l+s+s2}{\PYZdq{}}\PY{p}{,}
            \PY{l+s+sa}{f}\PY{l+s+s2}{\PYZdq{}}\PY{l+s+si}{\PYZob{}}\PY{n}{p\PYZus{}value\PYZus{}v}\PY{l+s+si}{:}\PY{l+s+s2}{.4f}\PY{l+s+si}{\PYZcb{}}\PY{l+s+s2}{\PYZdq{}}\PY{p}{,}
            \PY{l+s+sa}{f}\PY{l+s+s2}{\PYZdq{}}\PY{l+s+si}{\PYZob{}}\PY{n}{alpha\PYZus{}v}\PY{l+s+si}{\PYZcb{}}\PY{l+s+s2}{\PYZdq{}}\PY{p}{,}
            \PY{l+s+sa}{f}\PY{l+s+s2}{\PYZdq{}}\PY{l+s+si}{\PYZob{}}\PY{n}{type\PYZus{}i\PYZus{}error}\PY{l+s+si}{\PYZcb{}}\PY{l+s+s2}{\PYZdq{}}\PY{p}{,}
        \PY{p}{]}\PY{p}{,}
    \PY{p}{\PYZcb{}}
\PY{p}{)}
\end{Verbatim}
\end{tcolorbox}

    \begin{tcolorbox}[breakable, size=fbox, boxrule=1pt, pad at break*=1mm,colback=cellbackground, colframe=cellborder]
\prompt{In}{incolor}{32}{\boxspacing}
\begin{Verbatim}[commandchars=\\\{\}]
\PY{n}{mo}\PY{o}{.}\PY{n}{md}\PY{p}{(}\PY{l+s+sa}{rf}\PY{l+s+s2}{\PYZdq{}\PYZdq{}\PYZdq{}}
\PY{l+s+s2}{\PYZsh{}\PYZsh{} Test d}\PY{l+s+s2}{\PYZsq{}}\PY{l+s+s2}{hypothèse unilatéral à gauche}

\PY{l+s+s2}{**Formulation des hypothèses:**}
\PY{l+s+s2}{\PYZhy{} H₀: μ ≥ }\PY{l+s+si}{\PYZob{}}\PY{n}{mu\PYZus{}0}\PY{l+s+si}{\PYZcb{}}\PY{l+s+s2}{ minutes (}\PY{l+s+si}{\PYZob{}}\PY{n}{hours}\PY{p}{(}\PY{n}{mu\PYZus{}0}\PY{p}{)}\PY{l+s+si}{:}\PY{l+s+s2}{.2f}\PY{l+s+si}{\PYZcb{}}\PY{l+s+s2}{ heures) \PYZhy{} L}\PY{l+s+s2}{\PYZsq{}}\PY{l+s+s2}{hypothèse du patron}
\PY{l+s+s2}{\PYZhy{} H₁: μ \PYZlt{} }\PY{l+s+si}{\PYZob{}}\PY{n}{mu\PYZus{}0}\PY{l+s+si}{\PYZcb{}}\PY{l+s+s2}{ minutes \PYZhy{} Le temps de jeu moyen est inférieur à 5 heures}

\PY{l+s+s2}{**Méthode:** Test Z (n ≥ 30, distribution normale)}

\PY{l+s+s2}{La statistique de test est:}
\PY{l+s+s2}{\PYZdl{}\PYZdl{}}
\PY{l+s+s2}{z = }\PY{l+s+s2}{\PYZbs{}}\PY{l+s+s2}{frac}\PY{l+s+se}{\PYZob{}\PYZob{}}\PY{l+s+s2}{\PYZbs{}}\PY{l+s+s2}{bar}\PY{l+s+se}{\PYZob{}\PYZob{}}\PY{l+s+s2}{x}\PY{l+s+se}{\PYZcb{}\PYZcb{}}\PY{l+s+s2}{ \PYZhy{} }\PY{l+s+s2}{\PYZbs{}}\PY{l+s+s2}{mu\PYZus{}0}\PY{l+s+se}{\PYZcb{}\PYZcb{}}\PY{l+s+se}{\PYZob{}\PYZob{}}\PY{l+s+s2}{s / }\PY{l+s+s2}{\PYZbs{}}\PY{l+s+s2}{sqrt}\PY{l+s+se}{\PYZob{}\PYZob{}}\PY{l+s+s2}{n}\PY{l+s+se}{\PYZcb{}\PYZcb{}}\PY{l+s+se}{\PYZcb{}\PYZcb{}}\PY{l+s+s2}{ = }\PY{l+s+si}{\PYZob{}}\PY{n}{z\PYZus{}stat}\PY{l+s+si}{:}\PY{l+s+s2}{.4f}\PY{l+s+si}{\PYZcb{}}
\PY{l+s+s2}{\PYZdl{}\PYZdl{}}

\PY{l+s+s2}{**Résultats:**}

\PY{l+s+si}{\PYZob{}}\PY{n}{mo}\PY{o}{.}\PY{n}{ui}\PY{o}{.}\PY{n}{table}\PY{p}{(}\PY{n}{hyp\PYZus{}test\PYZus{}df}\PY{p}{,}\PY{+w}{ }\PY{n}{selection}\PY{o}{=}\PY{k+kc}{None}\PY{p}{)}\PY{l+s+si}{\PYZcb{}}

\PY{l+s+s2}{**Règle de décision:**}

\PY{l+s+s2}{On rejette H₀ si z \PYZlt{} z\PYZus{}α = }\PY{l+s+si}{\PYZob{}}\PY{n}{z\PYZus{}critical\PYZus{}v}\PY{l+s+si}{:}\PY{l+s+s2}{.4f}\PY{l+s+si}{\PYZcb{}}

\PY{l+s+s2}{Puisque z = }\PY{l+s+si}{\PYZob{}}\PY{n}{z\PYZus{}stat}\PY{l+s+si}{:}\PY{l+s+s2}{.4f}\PY{l+s+si}{\PYZcb{}}\PY{l+s+s2}{ }\PY{l+s+si}{\PYZob{}}\PY{l+s+s2}{\PYZdq{}}\PY{l+s+s2}{\PYZlt{}}\PY{l+s+s2}{\PYZdq{}}\PY{+w}{ }\PY{k}{if}\PY{+w}{ }\PY{n}{z\PYZus{}stat}\PY{+w}{ }\PY{o}{\PYZlt{}}\PY{+w}{ }\PY{n}{z\PYZus{}critical\PYZus{}v}\PY{+w}{ }\PY{k}{else}\PY{+w}{ }\PY{l+s+s2}{\PYZdq{}}\PY{l+s+s2}{≥}\PY{l+s+s2}{\PYZdq{}}\PY{l+s+si}{\PYZcb{}}\PY{l+s+s2}{ }\PY{l+s+si}{\PYZob{}}\PY{n}{z\PYZus{}critical\PYZus{}v}\PY{l+s+si}{:}\PY{l+s+s2}{.4f}\PY{l+s+si}{\PYZcb{}}\PY{l+s+s2}{, on }\PY{l+s+si}{\PYZob{}}\PY{l+s+s2}{\PYZdq{}}\PY{l+s+s2}{**rejette**}\PY{l+s+s2}{\PYZdq{}}\PY{+w}{ }\PY{k}{if}\PY{+w}{ }\PY{n}{reject\PYZus{}h0\PYZus{}v}\PY{+w}{ }\PY{k}{else}\PY{+w}{ }\PY{l+s+s2}{\PYZdq{}}\PY{l+s+s2}{**ne rejette pas**}\PY{l+s+s2}{\PYZdq{}}\PY{l+s+si}{\PYZcb{}}\PY{l+s+s2}{ H₀.}

\PY{l+s+s2}{**Conclusion:**}

\PY{l+s+s2}{Au niveau de signification α = }\PY{l+s+si}{\PYZob{}}\PY{n}{alpha\PYZus{}v}\PY{l+s+si}{\PYZcb{}}\PY{l+s+s2}{, }\PY{l+s+si}{\PYZob{}}\PY{l+s+s2}{\PYZdq{}}\PY{l+s+s2}{les données fournissent suffisamment de preuves pour rejeter l}\PY{l+s+s2}{\PYZsq{}}\PY{l+s+s2}{hypothèse du patron. Le temps de jeu moyen est significativement inférieur à 5 heures par semaine.}\PY{l+s+s2}{\PYZdq{}}\PY{+w}{ }\PY{k}{if}\PY{+w}{ }\PY{n}{reject\PYZus{}h0\PYZus{}v}\PY{+w}{ }\PY{k}{else}\PY{+w}{ }\PY{l+s+s2}{\PYZdq{}}\PY{l+s+s2}{les données ne fournissent pas suffisamment de preuves pour rejeter l}\PY{l+s+s2}{\PYZsq{}}\PY{l+s+s2}{hypothèse du patron. On ne peut pas conclure que le temps de jeu moyen est inférieur à 5 heures par semaine.}\PY{l+s+s2}{\PYZdq{}}\PY{l+s+si}{\PYZcb{}}

\PY{l+s+s2}{**Erreur de première espèce (α):**}

\PY{l+s+s2}{L}\PY{l+s+s2}{\PYZsq{}}\PY{l+s+s2}{erreur de type I est la probabilité de rejeter H₀ alors qu}\PY{l+s+s2}{\PYZsq{}}\PY{l+s+s2}{elle est vraie. Ici, α = }\PY{l+s+si}{\PYZob{}}\PY{n}{type\PYZus{}i\PYZus{}error}\PY{l+s+si}{\PYZcb{}}\PY{l+s+s2}{ = }\PY{l+s+si}{\PYZob{}}\PY{n}{type\PYZus{}i\PYZus{}error}\PY{+w}{ }\PY{o}{*}\PY{+w}{ }\PY{l+m+mi}{100}\PY{l+s+si}{\PYZcb{}}\PY{l+s+s2}{\PYZpc{}.}

\PY{l+s+s2}{Cela signifie qu}\PY{l+s+s2}{\PYZsq{}}\PY{l+s+s2}{il y a }\PY{l+s+si}{\PYZob{}}\PY{n}{type\PYZus{}i\PYZus{}error}\PY{+w}{ }\PY{o}{*}\PY{+w}{ }\PY{l+m+mi}{100}\PY{l+s+si}{\PYZcb{}}\PY{l+s+s2}{\PYZpc{} de chance de conclure à tort que le temps de jeu moyen est inférieur à 5 heures, alors qu}\PY{l+s+s2}{\PYZsq{}}\PY{l+s+s2}{en réalité il est d}\PY{l+s+s2}{\PYZsq{}}\PY{l+s+s2}{au moins 5 heures.}
\PY{l+s+s2}{\PYZdq{}\PYZdq{}\PYZdq{}}\PY{p}{)}
\end{Verbatim}
\end{tcolorbox}
 
            
\prompt{Out}{outcolor}{32}{}
    
    \subsection{Test d'hypothèse unilatéral à
gauche}\label{test-dhypothuxe8se-unilatuxe9ral-uxe0-gauche}

\textbf{Formulation des hypothèses:} - H₀: μ ≥ 300 minutes (5.00 heures)
- L'hypothèse du patron - H₁: μ \textless{} 300 minutes - Le temps de
jeu moyen est inférieur à 5 heures

\textbf{Méthode:} Test Z (n ≥ 30, distribution normale)

La statistique de test est: \[
z = \frac{\bar{x} - \mu_0}{s / \sqrt{n}} = -3.8549
\]

\textbf{Résultats:}

\textbf{Règle de décision:}

On rejette H₀ si z \textless{} z\_α = -1.6449

Puisque z = -3.8549 \textless{} -1.6449, on \textbf{rejette} H₀.

\textbf{Conclusion:}

Au niveau de signification α = 0.05, les données fournissent
suffisamment de preuves pour rejeter l'hypothèse du patron. Le temps de
jeu moyen est significativement inférieur à 5 heures par semaine.

\textbf{Erreur de première espèce (α):}

L'erreur de type I est la probabilité de rejeter H₀ alors qu'elle est
vraie. Ici, α = 0.05 = 5.0\%.

Cela signifie qu'il y a 5.0\% de chance de conclure à tort que le temps
de jeu moyen est inférieur à 5 heures, alors qu'en réalité il est d'au
moins 5 heures.

    

    \section{VI) Erreur de deuxième
espèce}\label{vi-erreur-de-deuxiuxe8me-espuxe8ce}

Cette section calcule l'erreur de type II (β) en supposant que la
moyenne échantillonnale est un bon estimé de la vraie moyenne de la
population, comme demandé dans la sixième partie du mandat.

Supposons que la moyenne échantillonnale est un très bon estimé de la
moyenne de la population et qu'elle peut être considérée comme étant
cette dernière, quelle est l'erreur de deuxième espèce commise au point
précédent si on ne rejette pas l'hypothèse nulle?

    \begin{tcolorbox}[breakable, size=fbox, boxrule=1pt, pad at break*=1mm,colback=cellbackground, colframe=cellborder]
\prompt{In}{incolor}{33}{\boxspacing}
\begin{Verbatim}[commandchars=\\\{\}]
\PY{c+c1}{\PYZsh{} Type II error (β) calculation}
\PY{c+c1}{\PYZsh{} β = P(not rejecting H0 | H1 is true)}
\PY{c+c1}{\PYZsh{} Assuming the true population mean is equal to the sample mean (x\PYZus{}bar)}

\PY{c+c1}{\PYZsh{} True population mean (assumed to be the sample mean)}
\PY{n}{mu\PYZus{}true} \PY{o}{=} \PY{n}{x\PYZus{}bar}

\PY{c+c1}{\PYZsh{} Standard error}
\PY{n}{se} \PY{o}{=} \PY{n}{s} \PY{o}{/} \PY{n}{np}\PY{o}{.}\PY{n}{sqrt}\PY{p}{(}\PY{n}{n}\PY{p}{)}

\PY{c+c1}{\PYZsh{} Critical value for the test (x\PYZus{}c such that we reject H0 if x̄ \PYZlt{} x\PYZus{}c)}
\PY{c+c1}{\PYZsh{} From z\PYZus{}α = (x\PYZus{}c \PYZhy{} μ₀) / se, we get x\PYZus{}c = μ₀ + z\PYZus{}α * se}
\PY{n}{z\PYZus{}alpha} \PY{o}{=} \PY{n}{stats}\PY{o}{.}\PY{n}{norm}\PY{o}{.}\PY{n}{ppf}\PY{p}{(}\PY{n}{alpha\PYZus{}v}\PY{p}{)}
\PY{n}{x\PYZus{}critical} \PY{o}{=} \PY{n}{mu\PYZus{}0} \PY{o}{+} \PY{n}{z\PYZus{}alpha} \PY{o}{*} \PY{n}{se}

\PY{c+c1}{\PYZsh{} Type II error: P(x̄ ≥ x\PYZus{}c | μ = μ\PYZus{}true)}
\PY{c+c1}{\PYZsh{} β = P(Z ≥ (x\PYZus{}c \PYZhy{} μ\PYZus{}true) / se)}
\PY{n}{z\PYZus{}beta} \PY{o}{=} \PY{p}{(}\PY{n}{x\PYZus{}critical} \PY{o}{\PYZhy{}} \PY{n}{mu\PYZus{}true}\PY{p}{)} \PY{o}{/} \PY{n}{se}
\PY{n}{beta} \PY{o}{=} \PY{l+m+mi}{1} \PY{o}{\PYZhy{}} \PY{n}{stats}\PY{o}{.}\PY{n}{norm}\PY{o}{.}\PY{n}{cdf}\PY{p}{(}\PY{n}{z\PYZus{}beta}\PY{p}{)}

\PY{c+c1}{\PYZsh{} Power of the test}
\PY{n}{power} \PY{o}{=} \PY{l+m+mi}{1} \PY{o}{\PYZhy{}} \PY{n}{beta}

\PY{n}{type\PYZus{}ii\PYZus{}df} \PY{o}{=} \PY{n}{pl}\PY{o}{.}\PY{n}{DataFrame}\PY{p}{(}
    \PY{p}{\PYZob{}}
        \PY{l+s+s2}{\PYZdq{}}\PY{l+s+s2}{Paramètre}\PY{l+s+s2}{\PYZdq{}}\PY{p}{:} \PY{p}{[}
            \PY{l+s+s2}{\PYZdq{}}\PY{l+s+s2}{Vraie moyenne de la population μ (min)}\PY{l+s+s2}{\PYZdq{}}\PY{p}{,}
            \PY{l+s+s2}{\PYZdq{}}\PY{l+s+s2}{Moyenne hypothétisée μ₀ (min)}\PY{l+s+s2}{\PYZdq{}}\PY{p}{,}
            \PY{l+s+s2}{\PYZdq{}}\PY{l+s+s2}{Erreur standard}\PY{l+s+s2}{\PYZdq{}}\PY{p}{,}
            \PY{l+s+s2}{\PYZdq{}}\PY{l+s+s2}{Valeur critique x\PYZus{}c (min)}\PY{l+s+s2}{\PYZdq{}}\PY{p}{,}
            \PY{l+s+s2}{\PYZdq{}}\PY{l+s+s2}{z\PYZus{}β}\PY{l+s+s2}{\PYZdq{}}\PY{p}{,}
            \PY{l+s+s2}{\PYZdq{}}\PY{l+s+s2}{Erreur de type II (β)}\PY{l+s+s2}{\PYZdq{}}\PY{p}{,}
            \PY{l+s+s2}{\PYZdq{}}\PY{l+s+s2}{Puissance du test (1 \PYZhy{} β)}\PY{l+s+s2}{\PYZdq{}}\PY{p}{,}
        \PY{p}{]}\PY{p}{,}
        \PY{l+s+s2}{\PYZdq{}}\PY{l+s+s2}{Valeur}\PY{l+s+s2}{\PYZdq{}}\PY{p}{:} \PY{p}{[}
            \PY{l+s+sa}{f}\PY{l+s+s2}{\PYZdq{}}\PY{l+s+si}{\PYZob{}}\PY{n}{mu\PYZus{}true}\PY{l+s+si}{:}\PY{l+s+s2}{.2f}\PY{l+s+si}{\PYZcb{}}\PY{l+s+s2}{\PYZdq{}}\PY{p}{,}
            \PY{l+s+sa}{f}\PY{l+s+s2}{\PYZdq{}}\PY{l+s+si}{\PYZob{}}\PY{n}{mu\PYZus{}0}\PY{l+s+si}{\PYZcb{}}\PY{l+s+s2}{\PYZdq{}}\PY{p}{,}
            \PY{l+s+sa}{f}\PY{l+s+s2}{\PYZdq{}}\PY{l+s+si}{\PYZob{}}\PY{n}{se}\PY{l+s+si}{:}\PY{l+s+s2}{.4f}\PY{l+s+si}{\PYZcb{}}\PY{l+s+s2}{\PYZdq{}}\PY{p}{,}
            \PY{l+s+sa}{f}\PY{l+s+s2}{\PYZdq{}}\PY{l+s+si}{\PYZob{}}\PY{n}{x\PYZus{}critical}\PY{l+s+si}{:}\PY{l+s+s2}{.2f}\PY{l+s+si}{\PYZcb{}}\PY{l+s+s2}{\PYZdq{}}\PY{p}{,}
            \PY{l+s+sa}{f}\PY{l+s+s2}{\PYZdq{}}\PY{l+s+si}{\PYZob{}}\PY{n}{z\PYZus{}beta}\PY{l+s+si}{:}\PY{l+s+s2}{.4f}\PY{l+s+si}{\PYZcb{}}\PY{l+s+s2}{\PYZdq{}}\PY{p}{,}
            \PY{l+s+sa}{f}\PY{l+s+s2}{\PYZdq{}}\PY{l+s+si}{\PYZob{}}\PY{n}{beta}\PY{l+s+si}{:}\PY{l+s+s2}{.4f}\PY{l+s+si}{\PYZcb{}}\PY{l+s+s2}{\PYZdq{}}\PY{p}{,}
            \PY{l+s+sa}{f}\PY{l+s+s2}{\PYZdq{}}\PY{l+s+si}{\PYZob{}}\PY{n}{power}\PY{l+s+si}{:}\PY{l+s+s2}{.4f}\PY{l+s+si}{\PYZcb{}}\PY{l+s+s2}{\PYZdq{}}\PY{p}{,}
        \PY{p}{]}\PY{p}{,}
    \PY{p}{\PYZcb{}}
\PY{p}{)}
\end{Verbatim}
\end{tcolorbox}

    \begin{tcolorbox}[breakable, size=fbox, boxrule=1pt, pad at break*=1mm,colback=cellbackground, colframe=cellborder]
\prompt{In}{incolor}{34}{\boxspacing}
\begin{Verbatim}[commandchars=\\\{\}]
\PY{n}{mo}\PY{o}{.}\PY{n}{md}\PY{p}{(}\PY{l+s+sa}{rf}\PY{l+s+s2}{\PYZdq{}\PYZdq{}\PYZdq{}}
\PY{l+s+s2}{\PYZsh{}\PYZsh{} Calcul de l}\PY{l+s+s2}{\PYZsq{}}\PY{l+s+s2}{erreur de deuxième espèce (β)}

\PY{l+s+s2}{L}\PY{l+s+s2}{\PYZsq{}}\PY{l+s+s2}{erreur de type II (β) est la probabilité de ne pas rejeter H₀ alors que H₁ est vraie.}

\PY{l+s+s2}{**Hypothèses:**}
\PY{l+s+s2}{\PYZhy{} On suppose que la vraie moyenne de la population est μ = x̄ = }\PY{l+s+si}{\PYZob{}}\PY{n}{mu\PYZus{}true}\PY{l+s+si}{:}\PY{l+s+s2}{.2f}\PY{l+s+si}{\PYZcb{}}\PY{l+s+s2}{ minutes (}\PY{l+s+si}{\PYZob{}}\PY{n}{hours}\PY{p}{(}\PY{n}{mu\PYZus{}true}\PY{p}{)}\PY{l+s+si}{:}\PY{l+s+s2}{.2f}\PY{l+s+si}{\PYZcb{}}\PY{l+s+s2}{ heures)}
\PY{l+s+s2}{\PYZhy{} H₀: μ ≥ }\PY{l+s+si}{\PYZob{}}\PY{n}{mu\PYZus{}0}\PY{l+s+si}{\PYZcb{}}\PY{l+s+s2}{ minutes}
\PY{l+s+s2}{\PYZhy{} H₁: μ \PYZlt{} }\PY{l+s+si}{\PYZob{}}\PY{n}{mu\PYZus{}0}\PY{l+s+si}{\PYZcb{}}\PY{l+s+s2}{ minutes}

\PY{l+s+s2}{**Calcul:**}

\PY{l+s+s2}{1. La valeur critique x\PYZus{}c est déterminée par:}
\PY{l+s+s2}{\PYZdl{}\PYZdl{}}
\PY{l+s+s2}{x\PYZus{}c = }\PY{l+s+s2}{\PYZbs{}}\PY{l+s+s2}{mu\PYZus{}0 + z\PYZus{}}\PY{l+s+s2}{\PYZbs{}}\PY{l+s+s2}{alpha }\PY{l+s+s2}{\PYZbs{}}\PY{l+s+s2}{times SE = }\PY{l+s+si}{\PYZob{}}\PY{n}{mu\PYZus{}0}\PY{l+s+si}{\PYZcb{}}\PY{l+s+s2}{ + }\PY{l+s+si}{\PYZob{}}\PY{n}{stats}\PY{o}{.}\PY{n}{norm}\PY{o}{.}\PY{n}{ppf}\PY{p}{(}\PY{n}{alpha\PYZus{}v}\PY{p}{)}\PY{l+s+si}{:}\PY{l+s+s2}{.4f}\PY{l+s+si}{\PYZcb{}}\PY{l+s+s2}{ }\PY{l+s+s2}{\PYZbs{}}\PY{l+s+s2}{times }\PY{l+s+si}{\PYZob{}}\PY{n}{se}\PY{l+s+si}{:}\PY{l+s+s2}{.4f}\PY{l+s+si}{\PYZcb{}}\PY{l+s+s2}{ = }\PY{l+s+si}{\PYZob{}}\PY{n}{x\PYZus{}critical}\PY{l+s+si}{:}\PY{l+s+s2}{.2f}\PY{l+s+si}{\PYZcb{}}
\PY{l+s+s2}{\PYZdl{}\PYZdl{}}

\PY{l+s+s2}{2. L}\PY{l+s+s2}{\PYZsq{}}\PY{l+s+s2}{erreur de type II est:}
\PY{l+s+s2}{\PYZdl{}\PYZdl{}}
\PY{l+s+s2}{\PYZbs{}}\PY{l+s+s2}{beta = P(}\PY{l+s+s2}{\PYZbs{}}\PY{l+s+s2}{bar}\PY{l+s+se}{\PYZob{}\PYZob{}}\PY{l+s+s2}{X}\PY{l+s+se}{\PYZcb{}\PYZcb{}}\PY{l+s+s2}{ }\PY{l+s+s2}{\PYZbs{}}\PY{l+s+s2}{geq x\PYZus{}c | }\PY{l+s+s2}{\PYZbs{}}\PY{l+s+s2}{mu = }\PY{l+s+s2}{\PYZbs{}}\PY{l+s+s2}{mu\PYZus{}}\PY{l+s+se}{\PYZob{}\PYZob{}}\PY{l+s+s2}{vrai}\PY{l+s+se}{\PYZcb{}\PYZcb{}}\PY{l+s+s2}{) = P}\PY{l+s+s2}{\PYZbs{}}\PY{l+s+s2}{left(Z }\PY{l+s+s2}{\PYZbs{}}\PY{l+s+s2}{geq }\PY{l+s+s2}{\PYZbs{}}\PY{l+s+s2}{frac}\PY{l+s+se}{\PYZob{}\PYZob{}}\PY{l+s+s2}{x\PYZus{}c \PYZhy{} }\PY{l+s+s2}{\PYZbs{}}\PY{l+s+s2}{mu\PYZus{}}\PY{l+s+se}{\PYZob{}\PYZob{}}\PY{l+s+s2}{vrai}\PY{l+s+se}{\PYZcb{}\PYZcb{}}\PY{l+s+se}{\PYZcb{}\PYZcb{}}\PY{l+s+se}{\PYZob{}\PYZob{}}\PY{l+s+s2}{SE}\PY{l+s+se}{\PYZcb{}\PYZcb{}}\PY{l+s+s2}{\PYZbs{}}\PY{l+s+s2}{right) = P(Z }\PY{l+s+s2}{\PYZbs{}}\PY{l+s+s2}{geq }\PY{l+s+si}{\PYZob{}}\PY{n}{z\PYZus{}beta}\PY{l+s+si}{:}\PY{l+s+s2}{.4f}\PY{l+s+si}{\PYZcb{}}\PY{l+s+s2}{) = }\PY{l+s+si}{\PYZob{}}\PY{n}{beta}\PY{l+s+si}{:}\PY{l+s+s2}{.4f}\PY{l+s+si}{\PYZcb{}}
\PY{l+s+s2}{\PYZdl{}\PYZdl{}}

\PY{l+s+s2}{**Résultats:**}

\PY{l+s+si}{\PYZob{}}\PY{n}{mo}\PY{o}{.}\PY{n}{ui}\PY{o}{.}\PY{n}{table}\PY{p}{(}\PY{n}{type\PYZus{}ii\PYZus{}df}\PY{p}{,}\PY{+w}{ }\PY{n}{selection}\PY{o}{=}\PY{k+kc}{None}\PY{p}{)}\PY{l+s+si}{\PYZcb{}}

\PY{l+s+s2}{**Interprétation:**}

\PY{l+s+s2}{\PYZhy{} L}\PY{l+s+s2}{\PYZsq{}}\PY{l+s+s2}{erreur de type II (β) = }\PY{l+s+si}{\PYZob{}}\PY{n}{beta}\PY{l+s+si}{:}\PY{l+s+s2}{.4f}\PY{l+s+si}{\PYZcb{}}\PY{l+s+s2}{ = }\PY{l+s+si}{\PYZob{}}\PY{n}{beta}\PY{+w}{ }\PY{o}{*}\PY{+w}{ }\PY{l+m+mi}{100}\PY{l+s+si}{:}\PY{l+s+s2}{.2f}\PY{l+s+si}{\PYZcb{}}\PY{l+s+s2}{\PYZpc{}}
\PY{l+s+s2}{\PYZhy{} Cela signifie qu}\PY{l+s+s2}{\PYZsq{}}\PY{l+s+s2}{il y a }\PY{l+s+si}{\PYZob{}}\PY{n}{beta}\PY{+w}{ }\PY{o}{*}\PY{+w}{ }\PY{l+m+mi}{100}\PY{l+s+si}{:}\PY{l+s+s2}{.2f}\PY{l+s+si}{\PYZcb{}}\PY{l+s+s2}{\PYZpc{} de chance de ne pas rejeter l}\PY{l+s+s2}{\PYZsq{}}\PY{l+s+s2}{hypothèse du patron (μ ≥ 300 min) alors qu}\PY{l+s+s2}{\PYZsq{}}\PY{l+s+s2}{en réalité la vraie moyenne est de }\PY{l+s+si}{\PYZob{}}\PY{n}{mu\PYZus{}true}\PY{l+s+si}{:}\PY{l+s+s2}{.2f}\PY{l+s+si}{\PYZcb{}}\PY{l+s+s2}{ minutes.}
\PY{l+s+s2}{\PYZhy{} La puissance du test (1 \PYZhy{} β) = }\PY{l+s+si}{\PYZob{}}\PY{n}{power}\PY{l+s+si}{:}\PY{l+s+s2}{.4f}\PY{l+s+si}{\PYZcb{}}\PY{l+s+s2}{ = }\PY{l+s+si}{\PYZob{}}\PY{n}{power}\PY{+w}{ }\PY{o}{*}\PY{+w}{ }\PY{l+m+mi}{100}\PY{l+s+si}{:}\PY{l+s+s2}{.2f}\PY{l+s+si}{\PYZcb{}}\PY{l+s+s2}{\PYZpc{}}
\PY{l+s+s2}{\PYZdq{}\PYZdq{}\PYZdq{}}\PY{p}{)}
\end{Verbatim}
\end{tcolorbox}
 
            
\prompt{Out}{outcolor}{34}{}
    
    \subsection{Calcul de l'erreur de deuxième espèce
(β)}\label{calcul-de-lerreur-de-deuxiuxe8me-espuxe8ce-ux3b2}

L'erreur de type II (β) est la probabilité de ne pas rejeter H₀ alors
que H₁ est vraie.

\textbf{Hypothèses:} - On suppose que la vraie moyenne de la population
est μ = x̄ = 280.58 minutes (4.68 heures) - H₀: μ ≥ 300 minutes - H₁: μ
\textless{} 300 minutes

\textbf{Calcul:}

\begin{enumerate}
\def\labelenumi{\arabic{enumi}.}
\item
  La valeur critique x\_c est déterminée par: \[
  x_c = \mu_0 + z_\alpha \times SE = 300 + -1.6449 \times 5.0377 = 291.71
  \]
\item
  L'erreur de type II est: \[
  \beta = P(\bar{X} \geq x_c | \mu = \mu_{vrai}) = P\left(Z \geq \frac{x_c - \mu_{vrai}}{SE}\right) = P(Z \geq 2.2101) = 0.0135
  \]
\end{enumerate}

\textbf{Résultats:}

\textbf{Interprétation:}

\begin{itemize}
\tightlist
\item
  L'erreur de type II (β) = 0.0135 = 1.35\%
\item
  Cela signifie qu'il y a 1.35\% de chance de ne pas rejeter l'hypothèse
  du patron (μ ≥ 300 min) alors qu'en réalité la vraie moyenne est de
  280.58 minutes.
\item
  La puissance du test (1 - β) = 0.9865 = 98.65\%
\end{itemize}

    

    \section{VII) Test d'hypothèse bilatéral sur la
variance}\label{vii-test-dhypothuxe8se-bilatuxe9ral-sur-la-variance}

Cette section effectue un test du chi-deux bilatéral sur la variance
pour vérifier si l'écart-type des temps de jeu est significativement
différent de 50 minutes, comme exigé dans la septième partie du mandat.

En supposant que l'écart-type des temps de jeu est de 50, effectuez un
test d'hypothèse bilatéral sur la variance avec un seuil de
signification de 5\% afin d'évaluer si les données fournissent
suffisamment de preuves pour rejeter ou ne pas rejeter l'hypothèse
nulle.

    \begin{tcolorbox}[breakable, size=fbox, boxrule=1pt, pad at break*=1mm,colback=cellbackground, colframe=cellborder]
\prompt{In}{incolor}{35}{\boxspacing}
\begin{Verbatim}[commandchars=\\\{\}]
\PY{c+c1}{\PYZsh{} Bilateral hypothesis test on variance}
\PY{c+c1}{\PYZsh{} H0: σ² = 50² = 2500 (population variance equals hypothesized value)}
\PY{c+c1}{\PYZsh{} H1: σ² ≠ 2500 (population variance is different)}

\PY{n}{sigma\PYZus{}0} \PY{o}{=} \PY{l+m+mi}{50}  \PY{c+c1}{\PYZsh{} hypothesized population standard deviation}
\PY{n}{sigma\PYZus{}0\PYZus{}squared} \PY{o}{=} \PY{n}{sigma\PYZus{}0}\PY{o}{*}\PY{o}{*}\PY{l+m+mi}{2}  \PY{c+c1}{\PYZsh{} hypothesized population variance}
\PY{n}{alpha\PYZus{}vii} \PY{o}{=} \PY{l+m+mf}{0.05}  \PY{c+c1}{\PYZsh{} significance level}

\PY{c+c1}{\PYZsh{} Sample statistics}
\PY{n}{n\PYZus{}vii} \PY{o}{=} \PY{n+nb}{len}\PY{p}{(}\PY{n}{dataset}\PY{p}{)}
\PY{n}{s\PYZus{}squared\PYZus{}vii} \PY{o}{=} \PY{n}{dataset}\PY{o}{.}\PY{n}{var}\PY{p}{(}\PY{n}{ddof}\PY{o}{=}\PY{l+m+mi}{1}\PY{p}{)}  \PY{c+c1}{\PYZsh{} sample variance (unbiased)}
\PY{n}{s\PYZus{}vii} \PY{o}{=} \PY{n}{dataset}\PY{o}{.}\PY{n}{std}\PY{p}{(}\PY{n}{ddof}\PY{o}{=}\PY{l+m+mi}{1}\PY{p}{)}  \PY{c+c1}{\PYZsh{} sample standard deviation}

\PY{c+c1}{\PYZsh{} Test statistic: Chi\PYZhy{}square}
\PY{c+c1}{\PYZsh{} χ² = (n\PYZhy{}1) * s² / σ₀²}
\PY{n}{chi2\PYZus{}stat\PYZus{}vii} \PY{o}{=} \PY{p}{(}\PY{n}{n\PYZus{}vii} \PY{o}{\PYZhy{}} \PY{l+m+mi}{1}\PY{p}{)} \PY{o}{*} \PY{n}{s\PYZus{}squared\PYZus{}vii} \PY{o}{/} \PY{n}{sigma\PYZus{}0\PYZus{}squared}

\PY{c+c1}{\PYZsh{} Degrees of freedom}
\PY{n}{df\PYZus{}vii} \PY{o}{=} \PY{n}{n\PYZus{}vii} \PY{o}{\PYZhy{}} \PY{l+m+mi}{1}

\PY{c+c1}{\PYZsh{} Critical values for two\PYZhy{}tailed test}
\PY{n}{chi2\PYZus{}lower} \PY{o}{=} \PY{n}{stats}\PY{o}{.}\PY{n}{chi2}\PY{o}{.}\PY{n}{ppf}\PY{p}{(}\PY{n}{alpha\PYZus{}vii} \PY{o}{/} \PY{l+m+mi}{2}\PY{p}{,} \PY{n}{df}\PY{o}{=}\PY{n}{df\PYZus{}vii}\PY{p}{)}
\PY{n}{chi2\PYZus{}upper} \PY{o}{=} \PY{n}{stats}\PY{o}{.}\PY{n}{chi2}\PY{o}{.}\PY{n}{ppf}\PY{p}{(}\PY{l+m+mi}{1} \PY{o}{\PYZhy{}} \PY{n}{alpha\PYZus{}vii} \PY{o}{/} \PY{l+m+mi}{2}\PY{p}{,} \PY{n}{df}\PY{o}{=}\PY{n}{df\PYZus{}vii}\PY{p}{)}

\PY{c+c1}{\PYZsh{} P\PYZhy{}value for two\PYZhy{}tailed test}
\PY{c+c1}{\PYZsh{} P\PYZhy{}value = 2 * min(P(χ² \PYZlt{} χ²\PYZus{}stat), P(χ² \PYZgt{} χ²\PYZus{}stat))}
\PY{n}{p\PYZus{}lower} \PY{o}{=} \PY{n}{stats}\PY{o}{.}\PY{n}{chi2}\PY{o}{.}\PY{n}{cdf}\PY{p}{(}\PY{n}{chi2\PYZus{}stat\PYZus{}vii}\PY{p}{,} \PY{n}{df}\PY{o}{=}\PY{n}{df\PYZus{}vii}\PY{p}{)}
\PY{n}{p\PYZus{}upper} \PY{o}{=} \PY{l+m+mi}{1} \PY{o}{\PYZhy{}} \PY{n}{p\PYZus{}lower}
\PY{n}{p\PYZus{}value\PYZus{}vii} \PY{o}{=} \PY{l+m+mi}{2} \PY{o}{*} \PY{n+nb}{min}\PY{p}{(}\PY{n}{p\PYZus{}lower}\PY{p}{,} \PY{n}{p\PYZus{}upper}\PY{p}{)}

\PY{c+c1}{\PYZsh{} Decision}
\PY{n}{reject\PYZus{}h0\PYZus{}vii} \PY{o}{=} \PY{n}{chi2\PYZus{}stat\PYZus{}vii} \PY{o}{\PYZlt{}} \PY{n}{chi2\PYZus{}lower} \PY{o+ow}{or} \PY{n}{chi2\PYZus{}stat\PYZus{}vii} \PY{o}{\PYZgt{}} \PY{n}{chi2\PYZus{}upper}

\PY{n}{variance\PYZus{}test\PYZus{}df} \PY{o}{=} \PY{n}{pl}\PY{o}{.}\PY{n}{DataFrame}\PY{p}{(}
    \PY{p}{\PYZob{}}
        \PY{l+s+s2}{\PYZdq{}}\PY{l+s+s2}{Paramètre}\PY{l+s+s2}{\PYZdq{}}\PY{p}{:} \PY{p}{[}
            \PY{l+s+s2}{\PYZdq{}}\PY{l+s+s2}{Hypothèse nulle H₀}\PY{l+s+s2}{\PYZdq{}}\PY{p}{,}
            \PY{l+s+s2}{\PYZdq{}}\PY{l+s+s2}{Hypothèse alternative H₁}\PY{l+s+s2}{\PYZdq{}}\PY{p}{,}
            \PY{l+s+s2}{\PYZdq{}}\PY{l+s+s2}{Écart\PYZhy{}type hypothétisé σ₀}\PY{l+s+s2}{\PYZdq{}}\PY{p}{,}
            \PY{l+s+s2}{\PYZdq{}}\PY{l+s+s2}{Variance hypothétisée σ₀²}\PY{l+s+s2}{\PYZdq{}}\PY{p}{,}
            \PY{l+s+s2}{\PYZdq{}}\PY{l+s+s2}{Taille de l}\PY{l+s+s2}{\PYZsq{}}\PY{l+s+s2}{échantillon n}\PY{l+s+s2}{\PYZdq{}}\PY{p}{,}
            \PY{l+s+s2}{\PYZdq{}}\PY{l+s+s2}{Écart\PYZhy{}type de l}\PY{l+s+s2}{\PYZsq{}}\PY{l+s+s2}{échantillon s}\PY{l+s+s2}{\PYZdq{}}\PY{p}{,}
            \PY{l+s+s2}{\PYZdq{}}\PY{l+s+s2}{Variance de l}\PY{l+s+s2}{\PYZsq{}}\PY{l+s+s2}{échantillon s²}\PY{l+s+s2}{\PYZdq{}}\PY{p}{,}
            \PY{l+s+s2}{\PYZdq{}}\PY{l+s+s2}{Degrés de liberté (n\PYZhy{}1)}\PY{l+s+s2}{\PYZdq{}}\PY{p}{,}
            \PY{l+s+s2}{\PYZdq{}}\PY{l+s+s2}{Statistique χ²}\PY{l+s+s2}{\PYZdq{}}\PY{p}{,}
            \PY{l+s+s2}{\PYZdq{}}\PY{l+s+s2}{Valeur critique inférieure χ²\PYZus{}α/2}\PY{l+s+s2}{\PYZdq{}}\PY{p}{,}
            \PY{l+s+s2}{\PYZdq{}}\PY{l+s+s2}{Valeur critique supérieure χ²\PYZus{}1\PYZhy{}α/2}\PY{l+s+s2}{\PYZdq{}}\PY{p}{,}
            \PY{l+s+s2}{\PYZdq{}}\PY{l+s+s2}{Valeur p}\PY{l+s+s2}{\PYZdq{}}\PY{p}{,}
            \PY{l+s+s2}{\PYZdq{}}\PY{l+s+s2}{Niveau de signification α}\PY{l+s+s2}{\PYZdq{}}\PY{p}{,}
        \PY{p}{]}\PY{p}{,}
        \PY{l+s+s2}{\PYZdq{}}\PY{l+s+s2}{Valeur}\PY{l+s+s2}{\PYZdq{}}\PY{p}{:} \PY{p}{[}
            \PY{l+s+s2}{\PYZdq{}}\PY{l+s+s2}{σ² = 2500}\PY{l+s+s2}{\PYZdq{}}\PY{p}{,}
            \PY{l+s+s2}{\PYZdq{}}\PY{l+s+s2}{σ² ≠ 2500}\PY{l+s+s2}{\PYZdq{}}\PY{p}{,}
            \PY{l+s+sa}{f}\PY{l+s+s2}{\PYZdq{}}\PY{l+s+si}{\PYZob{}}\PY{n}{sigma\PYZus{}0}\PY{l+s+si}{\PYZcb{}}\PY{l+s+s2}{\PYZdq{}}\PY{p}{,}
            \PY{l+s+sa}{f}\PY{l+s+s2}{\PYZdq{}}\PY{l+s+si}{\PYZob{}}\PY{n}{sigma\PYZus{}0\PYZus{}squared}\PY{l+s+si}{\PYZcb{}}\PY{l+s+s2}{\PYZdq{}}\PY{p}{,}
            \PY{l+s+sa}{f}\PY{l+s+s2}{\PYZdq{}}\PY{l+s+si}{\PYZob{}}\PY{n}{n\PYZus{}vii}\PY{l+s+si}{\PYZcb{}}\PY{l+s+s2}{\PYZdq{}}\PY{p}{,}
            \PY{l+s+sa}{f}\PY{l+s+s2}{\PYZdq{}}\PY{l+s+si}{\PYZob{}}\PY{n}{s\PYZus{}vii}\PY{l+s+si}{:}\PY{l+s+s2}{.2f}\PY{l+s+si}{\PYZcb{}}\PY{l+s+s2}{\PYZdq{}}\PY{p}{,}
            \PY{l+s+sa}{f}\PY{l+s+s2}{\PYZdq{}}\PY{l+s+si}{\PYZob{}}\PY{n}{s\PYZus{}squared\PYZus{}vii}\PY{l+s+si}{:}\PY{l+s+s2}{.2f}\PY{l+s+si}{\PYZcb{}}\PY{l+s+s2}{\PYZdq{}}\PY{p}{,}
            \PY{l+s+sa}{f}\PY{l+s+s2}{\PYZdq{}}\PY{l+s+si}{\PYZob{}}\PY{n}{df\PYZus{}vii}\PY{l+s+si}{\PYZcb{}}\PY{l+s+s2}{\PYZdq{}}\PY{p}{,}
            \PY{l+s+sa}{f}\PY{l+s+s2}{\PYZdq{}}\PY{l+s+si}{\PYZob{}}\PY{n}{chi2\PYZus{}stat\PYZus{}vii}\PY{l+s+si}{:}\PY{l+s+s2}{.4f}\PY{l+s+si}{\PYZcb{}}\PY{l+s+s2}{\PYZdq{}}\PY{p}{,}
            \PY{l+s+sa}{f}\PY{l+s+s2}{\PYZdq{}}\PY{l+s+si}{\PYZob{}}\PY{n}{chi2\PYZus{}lower}\PY{l+s+si}{:}\PY{l+s+s2}{.4f}\PY{l+s+si}{\PYZcb{}}\PY{l+s+s2}{\PYZdq{}}\PY{p}{,}
            \PY{l+s+sa}{f}\PY{l+s+s2}{\PYZdq{}}\PY{l+s+si}{\PYZob{}}\PY{n}{chi2\PYZus{}upper}\PY{l+s+si}{:}\PY{l+s+s2}{.4f}\PY{l+s+si}{\PYZcb{}}\PY{l+s+s2}{\PYZdq{}}\PY{p}{,}
            \PY{l+s+sa}{f}\PY{l+s+s2}{\PYZdq{}}\PY{l+s+si}{\PYZob{}}\PY{n}{p\PYZus{}value\PYZus{}vii}\PY{l+s+si}{:}\PY{l+s+s2}{.4f}\PY{l+s+si}{\PYZcb{}}\PY{l+s+s2}{\PYZdq{}}\PY{p}{,}
            \PY{l+s+sa}{f}\PY{l+s+s2}{\PYZdq{}}\PY{l+s+si}{\PYZob{}}\PY{n}{alpha\PYZus{}vii}\PY{l+s+si}{\PYZcb{}}\PY{l+s+s2}{\PYZdq{}}\PY{p}{,}
        \PY{p}{]}\PY{p}{,}
    \PY{p}{\PYZcb{}}
\PY{p}{)}
\end{Verbatim}
\end{tcolorbox}

    \begin{tcolorbox}[breakable, size=fbox, boxrule=1pt, pad at break*=1mm,colback=cellbackground, colframe=cellborder]
\prompt{In}{incolor}{36}{\boxspacing}
\begin{Verbatim}[commandchars=\\\{\}]
\PY{n}{mo}\PY{o}{.}\PY{n}{md}\PY{p}{(}\PY{l+s+sa}{rf}\PY{l+s+s2}{\PYZdq{}\PYZdq{}\PYZdq{}}
\PY{l+s+s2}{\PYZsh{}\PYZsh{} Test bilatéral sur la variance (test du Chi\PYZhy{}deux)}

\PY{l+s+s2}{**Formulation des hypothèses:**}
\PY{l+s+s2}{\PYZhy{} H₀: σ² = }\PY{l+s+si}{\PYZob{}}\PY{n}{sigma\PYZus{}0\PYZus{}squared}\PY{l+s+si}{\PYZcb{}}\PY{l+s+s2}{ (σ = }\PY{l+s+si}{\PYZob{}}\PY{n}{sigma\PYZus{}0}\PY{l+s+si}{\PYZcb{}}\PY{l+s+s2}{)}
\PY{l+s+s2}{\PYZhy{} H₁: σ² ≠ }\PY{l+s+si}{\PYZob{}}\PY{n}{sigma\PYZus{}0\PYZus{}squared}\PY{l+s+si}{\PYZcb{}}

\PY{l+s+s2}{**Méthode:** Test du Chi\PYZhy{}deux pour la variance}

\PY{l+s+s2}{La statistique de test est:}
\PY{l+s+s2}{\PYZdl{}\PYZdl{}}
\PY{l+s+s2}{\PYZbs{}}\PY{l+s+s2}{chi\PYZca{}2 = }\PY{l+s+s2}{\PYZbs{}}\PY{l+s+s2}{frac}\PY{l+s+se}{\PYZob{}\PYZob{}}\PY{l+s+s2}{(n\PYZhy{}1) }\PY{l+s+s2}{\PYZbs{}}\PY{l+s+s2}{cdot s\PYZca{}2}\PY{l+s+se}{\PYZcb{}\PYZcb{}}\PY{l+s+se}{\PYZob{}\PYZob{}}\PY{l+s+s2}{\PYZbs{}}\PY{l+s+s2}{sigma\PYZus{}0\PYZca{}2}\PY{l+s+se}{\PYZcb{}\PYZcb{}}\PY{l+s+s2}{ = }\PY{l+s+s2}{\PYZbs{}}\PY{l+s+s2}{frac}\PY{l+s+se}{\PYZob{}\PYZob{}}\PY{l+s+si}{\PYZob{}}\PY{n}{df\PYZus{}vii}\PY{l+s+si}{\PYZcb{}}\PY{l+s+s2}{ }\PY{l+s+s2}{\PYZbs{}}\PY{l+s+s2}{cdot }\PY{l+s+si}{\PYZob{}}\PY{n}{s\PYZus{}squared\PYZus{}vii}\PY{l+s+si}{:}\PY{l+s+s2}{.2f}\PY{l+s+se}{\PYZcb{}\PYZcb{}}\PY{l+s+si}{\PYZcb{}}\PY{l+s+se}{\PYZob{}\PYZob{}}\PY{l+s+si}{\PYZob{}}\PY{n}{sigma\PYZus{}0\PYZus{}squared}\PY{l+s+si}{\PYZcb{}}\PY{l+s+se}{\PYZcb{}\PYZcb{}}\PY{l+s+s2}{ = }\PY{l+s+si}{\PYZob{}}\PY{n}{chi2\PYZus{}stat\PYZus{}vii}\PY{l+s+si}{:}\PY{l+s+s2}{.4f}\PY{l+s+si}{\PYZcb{}}
\PY{l+s+s2}{\PYZdl{}\PYZdl{}}

\PY{l+s+s2}{**Résultats:**}

\PY{l+s+si}{\PYZob{}}\PY{n}{mo}\PY{o}{.}\PY{n}{ui}\PY{o}{.}\PY{n}{table}\PY{p}{(}\PY{n}{variance\PYZus{}test\PYZus{}df}\PY{p}{,}\PY{+w}{ }\PY{n}{selection}\PY{o}{=}\PY{k+kc}{None}\PY{p}{)}\PY{l+s+si}{\PYZcb{}}

\PY{l+s+s2}{**Règle de décision:**}

\PY{l+s+s2}{On rejette H₀ si χ² \PYZlt{} χ²\PYZus{}α/2 = }\PY{l+s+si}{\PYZob{}}\PY{n}{chi2\PYZus{}lower}\PY{l+s+si}{:}\PY{l+s+s2}{.4f}\PY{l+s+si}{\PYZcb{}}\PY{l+s+s2}{ ou χ² \PYZgt{} χ²\PYZus{}1\PYZhy{}α/2 = }\PY{l+s+si}{\PYZob{}}\PY{n}{chi2\PYZus{}upper}\PY{l+s+si}{:}\PY{l+s+s2}{.4f}\PY{l+s+si}{\PYZcb{}}

\PY{l+s+s2}{Puisque χ² = }\PY{l+s+si}{\PYZob{}}\PY{n}{chi2\PYZus{}stat\PYZus{}vii}\PY{l+s+si}{:}\PY{l+s+s2}{.4f}\PY{l+s+si}{\PYZcb{}}\PY{l+s+s2}{ }\PY{l+s+si}{\PYZob{}}\PY{l+s+s2}{\PYZdq{}}\PY{l+s+s2}{est dans la région de rejet}\PY{l+s+s2}{\PYZdq{}}\PY{+w}{ }\PY{k}{if}\PY{+w}{ }\PY{n}{reject\PYZus{}h0\PYZus{}vii}\PY{+w}{ }\PY{k}{else}\PY{+w}{ }\PY{l+s+s2}{\PYZdq{}}\PY{l+s+s2}{n}\PY{l+s+s2}{\PYZsq{}}\PY{l+s+s2}{est pas dans la région de rejet}\PY{l+s+s2}{\PYZdq{}}\PY{l+s+si}{\PYZcb{}}\PY{l+s+s2}{ [}\PY{l+s+si}{\PYZob{}}\PY{n}{chi2\PYZus{}lower}\PY{l+s+si}{:}\PY{l+s+s2}{.4f}\PY{l+s+si}{\PYZcb{}}\PY{l+s+s2}{, }\PY{l+s+si}{\PYZob{}}\PY{n}{chi2\PYZus{}upper}\PY{l+s+si}{:}\PY{l+s+s2}{.4f}\PY{l+s+si}{\PYZcb{}}\PY{l+s+s2}{], on }\PY{l+s+si}{\PYZob{}}\PY{l+s+s2}{\PYZdq{}}\PY{l+s+s2}{**rejette**}\PY{l+s+s2}{\PYZdq{}}\PY{+w}{ }\PY{k}{if}\PY{+w}{ }\PY{n}{reject\PYZus{}h0\PYZus{}vii}\PY{+w}{ }\PY{k}{else}\PY{+w}{ }\PY{l+s+s2}{\PYZdq{}}\PY{l+s+s2}{**ne rejette pas**}\PY{l+s+s2}{\PYZdq{}}\PY{l+s+si}{\PYZcb{}}\PY{l+s+s2}{ H₀.}

\PY{l+s+s2}{**Conclusion:**}

\PY{l+s+s2}{Au niveau de signification α = }\PY{l+s+si}{\PYZob{}}\PY{n}{alpha\PYZus{}vii}\PY{l+s+si}{\PYZcb{}}\PY{l+s+s2}{, }\PY{l+s+si}{\PYZob{}}\PY{l+s+s2}{\PYZdq{}}\PY{l+s+s2}{les données fournissent suffisamment de preuves pour rejeter l}\PY{l+s+s2}{\PYZsq{}}\PY{l+s+s2}{hypothèse que l}\PY{l+s+s2}{\PYZsq{}}\PY{l+s+s2}{écart\PYZhy{}type de la population est de 50 minutes. La variance de la population est significativement différente de 2500.}\PY{l+s+s2}{\PYZdq{}}\PY{+w}{ }\PY{k}{if}\PY{+w}{ }\PY{n}{reject\PYZus{}h0\PYZus{}vii}\PY{+w}{ }\PY{k}{else}\PY{+w}{ }\PY{l+s+s2}{\PYZdq{}}\PY{l+s+s2}{les données ne fournissent pas suffisamment de preuves pour rejeter l}\PY{l+s+s2}{\PYZsq{}}\PY{l+s+s2}{hypothèse nulle. On ne peut pas conclure que la variance de la population est différente de 2500.}\PY{l+s+s2}{\PYZdq{}}\PY{l+s+si}{\PYZcb{}}
\PY{l+s+s2}{\PYZdq{}\PYZdq{}\PYZdq{}}\PY{p}{)}
\end{Verbatim}
\end{tcolorbox}
 
            
\prompt{Out}{outcolor}{36}{}
    
    \subsection{Test bilatéral sur la variance (test du
Chi-deux)}\label{test-bilatuxe9ral-sur-la-variance-test-du-chi-deux}

\textbf{Formulation des hypothèses:} - H₀: σ² = 2500 (σ = 50) - H₁: σ² ≠
2500

\textbf{Méthode:} Test du Chi-deux pour la variance

La statistique de test est: \[
\chi^2 = \frac{(n-1) \cdot s^2}{\sigma_0^2} = \frac{99 \cdot 2537.84}{2500} = 100.4985
\]

\textbf{Résultats:}

\textbf{Règle de décision:}

On rejette H₀ si χ² \textless{} χ²\_α/2 = 73.3611 ou χ² \textgreater{}
χ²\_1-α/2 = 128.4220

Puisque χ² = 100.4985 n'est pas dans la région de rejet {[}73.3611,
128.4220{]}, on \textbf{ne rejette pas} H₀.

\textbf{Conclusion:}

Au niveau de signification α = 0.05, les données ne fournissent pas
suffisamment de preuves pour rejeter l'hypothèse nulle. On ne peut pas
conclure que la variance de la population est différente de 2500.

    

    \section{Extra) Variable aléatoire Q du temps de
jeu}\label{extra-variable-aluxe9atoire-q-du-temps-de-jeu}

Cette section définit la variable aléatoire Q représentant le temps de
jeu hebdomadaire d'un joueur, qui sera utilisée dans le mandat 3 pour
les simulations Monte-Carlo.

On définit ici la variable aléatoire Q représentant le temps de jeu d'un
joueur.

    \begin{tcolorbox}[breakable, size=fbox, boxrule=1pt, pad at break*=1mm,colback=cellbackground, colframe=cellborder]
\prompt{In}{incolor}{37}{\boxspacing}
\begin{Verbatim}[commandchars=\\\{\}]
\PY{n}{mu\PYZus{}q} \PY{o}{=} \PY{n}{dataset}\PY{o}{.}\PY{n}{mean}\PY{p}{(}\PY{p}{)}
\PY{n}{sigma\PYZus{}q} \PY{o}{=} \PY{n}{dataset}\PY{o}{.}\PY{n}{std}\PY{p}{(}\PY{n}{ddof}\PY{o}{=}\PY{l+m+mi}{1}\PY{p}{)}
\end{Verbatim}
\end{tcolorbox}

    \begin{tcolorbox}[breakable, size=fbox, boxrule=1pt, pad at break*=1mm,colback=cellbackground, colframe=cellborder]
\prompt{In}{incolor}{38}{\boxspacing}
\begin{Verbatim}[commandchars=\\\{\}]
\PY{n}{mo}\PY{o}{.}\PY{n}{md}\PY{p}{(}\PY{l+s+sa}{rf}\PY{l+s+s2}{\PYZdq{}\PYZdq{}\PYZdq{}}
\PY{l+s+s2}{\PYZsh{}\PYZsh{} Définition de la variable aléatoire Q}

\PY{l+s+s2}{Soit **Q** la variable aléatoire représentant le temps de jeu hebdomadaire d}\PY{l+s+s2}{\PYZsq{}}\PY{l+s+s2}{un joueur (en minutes).}

\PY{l+s+s2}{En supposant que les temps de jeu suivent une distribution normale, on a :}

\PY{l+s+s2}{\PYZdl{}\PYZdl{}Q }\PY{l+s+s2}{\PYZbs{}}\PY{l+s+s2}{sim }\PY{l+s+s2}{\PYZbs{}}\PY{l+s+s2}{mathcal}\PY{l+s+se}{\PYZob{}\PYZob{}}\PY{l+s+s2}{N}\PY{l+s+se}{\PYZcb{}\PYZcb{}}\PY{l+s+s2}{(}\PY{l+s+s2}{\PYZbs{}}\PY{l+s+s2}{mu, }\PY{l+s+s2}{\PYZbs{}}\PY{l+s+s2}{sigma\PYZca{}2)\PYZdl{}\PYZdl{}}

\PY{l+s+s2}{Où :}

\PY{l+s+s2}{| Paramètre | Description | Valeur (min.) |}
\PY{l+s+s2}{|\PYZhy{}\PYZhy{}\PYZhy{}\PYZhy{}\PYZhy{}\PYZhy{}\PYZhy{}\PYZhy{}\PYZhy{}\PYZhy{}\PYZhy{}|\PYZhy{}\PYZhy{}\PYZhy{}\PYZhy{}\PYZhy{}\PYZhy{}\PYZhy{}\PYZhy{}\PYZhy{}\PYZhy{}\PYZhy{}\PYZhy{}\PYZhy{}|\PYZhy{}\PYZhy{}\PYZhy{}\PYZhy{}\PYZhy{}\PYZhy{}\PYZhy{}\PYZhy{}\PYZhy{}\PYZhy{}\PYZhy{}\PYZhy{}\PYZhy{}\PYZhy{}|}
\PY{l+s+s2}{| \PYZdl{}}\PY{l+s+s2}{\PYZbs{}}\PY{l+s+s2}{mu\PYZdl{} | Moyenne | }\PY{l+s+si}{\PYZob{}}\PY{n}{mu\PYZus{}q}\PY{l+s+si}{:}\PY{l+s+s2}{.2f}\PY{l+s+si}{\PYZcb{}}\PY{l+s+s2}{ |}
\PY{l+s+s2}{| \PYZdl{}}\PY{l+s+s2}{\PYZbs{}}\PY{l+s+s2}{sigma\PYZdl{} | Écart\PYZhy{}type | }\PY{l+s+si}{\PYZob{}}\PY{n}{sigma\PYZus{}q}\PY{l+s+si}{:}\PY{l+s+s2}{.2f}\PY{l+s+si}{\PYZcb{}}\PY{l+s+s2}{ |}
\PY{l+s+s2}{| \PYZdl{}}\PY{l+s+s2}{\PYZbs{}}\PY{l+s+s2}{sigma\PYZca{}2\PYZdl{} | Variance | }\PY{l+s+si}{\PYZob{}}\PY{n}{sigma\PYZus{}q}\PY{o}{*}\PY{o}{*}\PY{l+m+mi}{2}\PY{l+s+si}{:}\PY{l+s+s2}{.2f}\PY{l+s+si}{\PYZcb{}}\PY{l+s+s2}{ |}

\PY{l+s+s2}{Thus, the probability density function of Q is:}

\PY{l+s+s2}{\PYZdl{}\PYZdl{}f\PYZus{}Q(q) = }\PY{l+s+s2}{\PYZbs{}}\PY{l+s+s2}{frac}\PY{l+s+se}{\PYZob{}\PYZob{}}\PY{l+s+s2}{1}\PY{l+s+se}{\PYZcb{}\PYZcb{}}\PY{l+s+se}{\PYZob{}\PYZob{}}\PY{l+s+s2}{\PYZbs{}}\PY{l+s+s2}{sigma}\PY{l+s+s2}{\PYZbs{}}\PY{l+s+s2}{sqrt}\PY{l+s+se}{\PYZob{}\PYZob{}}\PY{l+s+s2}{2}\PY{l+s+s2}{\PYZbs{}}\PY{l+s+s2}{pi}\PY{l+s+se}{\PYZcb{}\PYZcb{}}\PY{l+s+se}{\PYZcb{}\PYZcb{}}\PY{l+s+s2}{ }\PY{l+s+s2}{\PYZbs{}}\PY{l+s+s2}{exp}\PY{l+s+s2}{\PYZbs{}}\PY{l+s+s2}{left(\PYZhy{}}\PY{l+s+s2}{\PYZbs{}}\PY{l+s+s2}{frac}\PY{l+s+se}{\PYZob{}\PYZob{}}\PY{l+s+s2}{(q\PYZhy{}}\PY{l+s+s2}{\PYZbs{}}\PY{l+s+s2}{mu)\PYZca{}2}\PY{l+s+se}{\PYZcb{}\PYZcb{}}\PY{l+s+se}{\PYZob{}\PYZob{}}\PY{l+s+s2}{2}\PY{l+s+s2}{\PYZbs{}}\PY{l+s+s2}{sigma\PYZca{}2}\PY{l+s+se}{\PYZcb{}\PYZcb{}}\PY{l+s+s2}{\PYZbs{}}\PY{l+s+s2}{right) = }\PY{l+s+s2}{\PYZbs{}}\PY{l+s+s2}{frac}\PY{l+s+se}{\PYZob{}\PYZob{}}\PY{l+s+s2}{1}\PY{l+s+se}{\PYZcb{}\PYZcb{}}\PY{l+s+se}{\PYZob{}\PYZob{}}\PY{l+s+si}{\PYZob{}}\PY{n}{sigma\PYZus{}q}\PY{l+s+si}{:}\PY{l+s+s2}{.2f}\PY{l+s+si}{\PYZcb{}}\PY{l+s+s2}{\PYZbs{}}\PY{l+s+s2}{sqrt}\PY{l+s+se}{\PYZob{}\PYZob{}}\PY{l+s+s2}{2}\PY{l+s+s2}{\PYZbs{}}\PY{l+s+s2}{pi}\PY{l+s+se}{\PYZcb{}\PYZcb{}}\PY{l+s+se}{\PYZcb{}\PYZcb{}}\PY{l+s+s2}{ }\PY{l+s+s2}{\PYZbs{}}\PY{l+s+s2}{exp}\PY{l+s+s2}{\PYZbs{}}\PY{l+s+s2}{left(\PYZhy{}}\PY{l+s+s2}{\PYZbs{}}\PY{l+s+s2}{frac}\PY{l+s+se}{\PYZob{}\PYZob{}}\PY{l+s+s2}{(q\PYZhy{}}\PY{l+s+si}{\PYZob{}}\PY{n}{mu\PYZus{}q}\PY{l+s+si}{:}\PY{l+s+s2}{.2f}\PY{l+s+si}{\PYZcb{}}\PY{l+s+s2}{)\PYZca{}2}\PY{l+s+se}{\PYZcb{}\PYZcb{}}\PY{l+s+se}{\PYZob{}\PYZob{}}\PY{l+s+s2}{2 }\PY{l+s+s2}{\PYZbs{}}\PY{l+s+s2}{times }\PY{l+s+si}{\PYZob{}}\PY{n}{sigma\PYZus{}q}\PY{o}{*}\PY{o}{*}\PY{l+m+mi}{2}\PY{l+s+si}{:}\PY{l+s+s2}{.2f}\PY{l+s+se}{\PYZcb{}\PYZcb{}}\PY{l+s+si}{\PYZcb{}}\PY{l+s+s2}{\PYZbs{}}\PY{l+s+s2}{right)\PYZdl{}\PYZdl{}}
\PY{l+s+s2}{\PYZdq{}\PYZdq{}\PYZdq{}}\PY{p}{)}
\end{Verbatim}
\end{tcolorbox}
 
            
\prompt{Out}{outcolor}{38}{}
    
    \subsection{Définition de la variable aléatoire
Q}\label{duxe9finition-de-la-variable-aluxe9atoire-q}

Soit \textbf{Q} la variable aléatoire représentant le temps de jeu
hebdomadaire d'un joueur (en minutes).

En supposant que les temps de jeu suivent une distribution normale, on a
:

\[Q \sim \mathcal{N}(\mu, \sigma^2)\]

Où :

\begin{longtable}[]{@{}lll@{}}
\toprule\noalign{}
Paramètre & Description & Valeur (min.) \\
\midrule\noalign{}
\endhead
\bottomrule\noalign{}
\endlastfoot
\(\mu\) & Moyenne & 280.58 \\
\(\sigma\) & Écart-type & 50.38 \\
\(\sigma^2\) & Variance & 2537.84 \\
\end{longtable}

Thus, the probability density function of Q is:

\[f_Q(q) = \frac{1}{\sigma\sqrt{2\pi}} \exp\left(-\frac{(q-\mu)^2}{2\sigma^2}\right) = \frac{1}{50.38\sqrt{2\pi}} \exp\left(-\frac{(q-280.58)^2}{2 \times 2537.84}\right)\]

    

    \begin{tcolorbox}[breakable, size=fbox, boxrule=1pt, pad at break*=1mm,colback=cellbackground, colframe=cellborder]
\prompt{In}{incolor}{ }{\boxspacing}
\begin{Verbatim}[commandchars=\\\{\}]

\end{Verbatim}
\end{tcolorbox}

    \begin{tcolorbox}[breakable, size=fbox, boxrule=1pt, pad at break*=1mm,colback=cellbackground, colframe=cellborder]
\prompt{In}{incolor}{ }{\boxspacing}
\begin{Verbatim}[commandchars=\\\{\}]

\end{Verbatim}
\end{tcolorbox}


    % Add a bibliography block to the postdoc
    
    
    
\end{document}
